\documentclass[11pt]{article}

    \usepackage[breakable]{tcolorbox}
    \usepackage{parskip} % Stop auto-indenting (to mimic markdown behaviour)
    

    % Basic figure setup, for now with no caption control since it's done
    % automatically by Pandoc (which extracts ![](path) syntax from Markdown).
    \usepackage{graphicx}
    % Keep aspect ratio if custom image width or height is specified
    \setkeys{Gin}{keepaspectratio}
    % Maintain compatibility with old templates. Remove in nbconvert 6.0
    \let\Oldincludegraphics\includegraphics
    % Ensure that by default, figures have no caption (until we provide a
    % proper Figure object with a Caption API and a way to capture that
    % in the conversion process - todo).
    \usepackage{caption}
    \DeclareCaptionFormat{nocaption}{}
    \captionsetup{format=nocaption,aboveskip=0pt,belowskip=0pt}

    \usepackage{float}
    \floatplacement{figure}{H} % forces figures to be placed at the correct location
    \usepackage{xcolor} % Allow colors to be defined
    \usepackage{enumerate} % Needed for markdown enumerations to work
    \usepackage{geometry} % Used to adjust the document margins
    \usepackage{amsmath} % Equations
    \usepackage{amssymb} % Equations
    \usepackage{textcomp} % defines textquotesingle
    % Hack from http://tex.stackexchange.com/a/47451/13684:
    \AtBeginDocument{%
        \def\PYZsq{\textquotesingle}% Upright quotes in Pygmentized code
    }
    \usepackage{upquote} % Upright quotes for verbatim code
    \usepackage{eurosym} % defines \euro

    \usepackage{iftex}
    \ifPDFTeX
        \usepackage[T1]{fontenc}
        \IfFileExists{alphabeta.sty}{
              \usepackage{alphabeta}
          }{
              \usepackage[mathletters]{ucs}
              \usepackage[utf8x]{inputenc}
          }
    \else
        \usepackage{fontspec}
        \usepackage{unicode-math}
    \fi

    \usepackage{fancyvrb} % verbatim replacement that allows latex
    \usepackage{grffile} % extends the file name processing of package graphics
                         % to support a larger range
    \makeatletter % fix for old versions of grffile with XeLaTeX
    \@ifpackagelater{grffile}{2019/11/01}
    {
      % Do nothing on new versions
    }
    {
      \def\Gread@@xetex#1{%
        \IfFileExists{"\Gin@base".bb}%
        {\Gread@eps{\Gin@base.bb}}%
        {\Gread@@xetex@aux#1}%
      }
    }
    \makeatother
    \usepackage[Export]{adjustbox} % Used to constrain images to a maximum size
    \adjustboxset{max size={0.9\linewidth}{0.9\paperheight}}

    % The hyperref package gives us a pdf with properly built
    % internal navigation ('pdf bookmarks' for the table of contents,
    % internal cross-reference links, web links for URLs, etc.)
    \usepackage{hyperref}
    % The default LaTeX title has an obnoxious amount of whitespace. By default,
    % titling removes some of it. It also provides customization options.
    \usepackage{titling}
    \usepackage{longtable} % longtable support required by pandoc >1.10
    \usepackage{booktabs}  % table support for pandoc > 1.12.2
    \usepackage{array}     % table support for pandoc >= 2.11.3
    \usepackage{calc}      % table minipage width calculation for pandoc >= 2.11.1
    \usepackage[inline]{enumitem} % IRkernel/repr support (it uses the enumerate* environment)
    \usepackage[normalem]{ulem} % ulem is needed to support strikethroughs (\sout)
                                % normalem makes italics be italics, not underlines
    \usepackage{soul}      % strikethrough (\st) support for pandoc >= 3.0.0
    \usepackage{mathrsfs}
    

    
    % Colors for the hyperref package
    \definecolor{urlcolor}{rgb}{0,.145,.698}
    \definecolor{linkcolor}{rgb}{.71,0.21,0.01}
    \definecolor{citecolor}{rgb}{.12,.54,.11}

    % ANSI colors
    \definecolor{ansi-black}{HTML}{3E424D}
    \definecolor{ansi-black-intense}{HTML}{282C36}
    \definecolor{ansi-red}{HTML}{E75C58}
    \definecolor{ansi-red-intense}{HTML}{B22B31}
    \definecolor{ansi-green}{HTML}{00A250}
    \definecolor{ansi-green-intense}{HTML}{007427}
    \definecolor{ansi-yellow}{HTML}{DDB62B}
    \definecolor{ansi-yellow-intense}{HTML}{B27D12}
    \definecolor{ansi-blue}{HTML}{208FFB}
    \definecolor{ansi-blue-intense}{HTML}{0065CA}
    \definecolor{ansi-magenta}{HTML}{D160C4}
    \definecolor{ansi-magenta-intense}{HTML}{A03196}
    \definecolor{ansi-cyan}{HTML}{60C6C8}
    \definecolor{ansi-cyan-intense}{HTML}{258F8F}
    \definecolor{ansi-white}{HTML}{C5C1B4}
    \definecolor{ansi-white-intense}{HTML}{A1A6B2}
    \definecolor{ansi-default-inverse-fg}{HTML}{FFFFFF}
    \definecolor{ansi-default-inverse-bg}{HTML}{000000}

    % common color for the border for error outputs.
    \definecolor{outerrorbackground}{HTML}{FFDFDF}

    % commands and environments needed by pandoc snippets
    % extracted from the output of `pandoc -s`
    \providecommand{\tightlist}{%
      \setlength{\itemsep}{0pt}\setlength{\parskip}{0pt}}
    \DefineVerbatimEnvironment{Highlighting}{Verbatim}{commandchars=\\\{\}}
    % Add ',fontsize=\small' for more characters per line
    \newenvironment{Shaded}{}{}
    \newcommand{\KeywordTok}[1]{\textcolor[rgb]{0.00,0.44,0.13}{\textbf{{#1}}}}
    \newcommand{\DataTypeTok}[1]{\textcolor[rgb]{0.56,0.13,0.00}{{#1}}}
    \newcommand{\DecValTok}[1]{\textcolor[rgb]{0.25,0.63,0.44}{{#1}}}
    \newcommand{\BaseNTok}[1]{\textcolor[rgb]{0.25,0.63,0.44}{{#1}}}
    \newcommand{\FloatTok}[1]{\textcolor[rgb]{0.25,0.63,0.44}{{#1}}}
    \newcommand{\CharTok}[1]{\textcolor[rgb]{0.25,0.44,0.63}{{#1}}}
    \newcommand{\StringTok}[1]{\textcolor[rgb]{0.25,0.44,0.63}{{#1}}}
    \newcommand{\CommentTok}[1]{\textcolor[rgb]{0.38,0.63,0.69}{\textit{{#1}}}}
    \newcommand{\OtherTok}[1]{\textcolor[rgb]{0.00,0.44,0.13}{{#1}}}
    \newcommand{\AlertTok}[1]{\textcolor[rgb]{1.00,0.00,0.00}{\textbf{{#1}}}}
    \newcommand{\FunctionTok}[1]{\textcolor[rgb]{0.02,0.16,0.49}{{#1}}}
    \newcommand{\RegionMarkerTok}[1]{{#1}}
    \newcommand{\ErrorTok}[1]{\textcolor[rgb]{1.00,0.00,0.00}{\textbf{{#1}}}}
    \newcommand{\NormalTok}[1]{{#1}}

    % Additional commands for more recent versions of Pandoc
    \newcommand{\ConstantTok}[1]{\textcolor[rgb]{0.53,0.00,0.00}{{#1}}}
    \newcommand{\SpecialCharTok}[1]{\textcolor[rgb]{0.25,0.44,0.63}{{#1}}}
    \newcommand{\VerbatimStringTok}[1]{\textcolor[rgb]{0.25,0.44,0.63}{{#1}}}
    \newcommand{\SpecialStringTok}[1]{\textcolor[rgb]{0.73,0.40,0.53}{{#1}}}
    \newcommand{\ImportTok}[1]{{#1}}
    \newcommand{\DocumentationTok}[1]{\textcolor[rgb]{0.73,0.13,0.13}{\textit{{#1}}}}
    \newcommand{\AnnotationTok}[1]{\textcolor[rgb]{0.38,0.63,0.69}{\textbf{\textit{{#1}}}}}
    \newcommand{\CommentVarTok}[1]{\textcolor[rgb]{0.38,0.63,0.69}{\textbf{\textit{{#1}}}}}
    \newcommand{\VariableTok}[1]{\textcolor[rgb]{0.10,0.09,0.49}{{#1}}}
    \newcommand{\ControlFlowTok}[1]{\textcolor[rgb]{0.00,0.44,0.13}{\textbf{{#1}}}}
    \newcommand{\OperatorTok}[1]{\textcolor[rgb]{0.40,0.40,0.40}{{#1}}}
    \newcommand{\BuiltInTok}[1]{{#1}}
    \newcommand{\ExtensionTok}[1]{{#1}}
    \newcommand{\PreprocessorTok}[1]{\textcolor[rgb]{0.74,0.48,0.00}{{#1}}}
    \newcommand{\AttributeTok}[1]{\textcolor[rgb]{0.49,0.56,0.16}{{#1}}}
    \newcommand{\InformationTok}[1]{\textcolor[rgb]{0.38,0.63,0.69}{\textbf{\textit{{#1}}}}}
    \newcommand{\WarningTok}[1]{\textcolor[rgb]{0.38,0.63,0.69}{\textbf{\textit{{#1}}}}}
    \makeatletter
    \newsavebox\pandoc@box
    \newcommand*\pandocbounded[1]{%
      \sbox\pandoc@box{#1}%
      % scaling factors for width and height
      \Gscale@div\@tempa\textheight{\dimexpr\ht\pandoc@box+\dp\pandoc@box\relax}%
      \Gscale@div\@tempb\linewidth{\wd\pandoc@box}%
      % select the smaller of both
      \ifdim\@tempb\p@<\@tempa\p@
        \let\@tempa\@tempb
      \fi
      % scaling accordingly (\@tempa < 1)
      \ifdim\@tempa\p@<\p@
        \scalebox{\@tempa}{\usebox\pandoc@box}%
      % scaling not needed, use as it is
      \else
        \usebox{\pandoc@box}%
      \fi
    }
    \makeatother

    % Define a nice break command that doesn't care if a line doesn't already
    % exist.
    \def\br{\hspace*{\fill} \\* }
    % Math Jax compatibility definitions
    \def\gt{>}
    \def\lt{<}
    \let\Oldtex\TeX
    \let\Oldlatex\LaTeX
    \renewcommand{\TeX}{\textrm{\Oldtex}}
    \renewcommand{\LaTeX}{\textrm{\Oldlatex}}
    % Document parameters
    % Document title
    \title{Notebook}
    
    
    
    
    
    
    
% Pygments definitions
\makeatletter
\def\PY@reset{\let\PY@it=\relax \let\PY@bf=\relax%
    \let\PY@ul=\relax \let\PY@tc=\relax%
    \let\PY@bc=\relax \let\PY@ff=\relax}
\def\PY@tok#1{\csname PY@tok@#1\endcsname}
\def\PY@toks#1+{\ifx\relax#1\empty\else%
    \PY@tok{#1}\expandafter\PY@toks\fi}
\def\PY@do#1{\PY@bc{\PY@tc{\PY@ul{%
    \PY@it{\PY@bf{\PY@ff{#1}}}}}}}
\def\PY#1#2{\PY@reset\PY@toks#1+\relax+\PY@do{#2}}

\@namedef{PY@tok@w}{\def\PY@tc##1{\textcolor[rgb]{0.73,0.73,0.73}{##1}}}
\@namedef{PY@tok@c}{\let\PY@it=\textit\def\PY@tc##1{\textcolor[rgb]{0.24,0.48,0.48}{##1}}}
\@namedef{PY@tok@cp}{\def\PY@tc##1{\textcolor[rgb]{0.61,0.40,0.00}{##1}}}
\@namedef{PY@tok@k}{\let\PY@bf=\textbf\def\PY@tc##1{\textcolor[rgb]{0.00,0.50,0.00}{##1}}}
\@namedef{PY@tok@kp}{\def\PY@tc##1{\textcolor[rgb]{0.00,0.50,0.00}{##1}}}
\@namedef{PY@tok@kt}{\def\PY@tc##1{\textcolor[rgb]{0.69,0.00,0.25}{##1}}}
\@namedef{PY@tok@o}{\def\PY@tc##1{\textcolor[rgb]{0.40,0.40,0.40}{##1}}}
\@namedef{PY@tok@ow}{\let\PY@bf=\textbf\def\PY@tc##1{\textcolor[rgb]{0.67,0.13,1.00}{##1}}}
\@namedef{PY@tok@nb}{\def\PY@tc##1{\textcolor[rgb]{0.00,0.50,0.00}{##1}}}
\@namedef{PY@tok@nf}{\def\PY@tc##1{\textcolor[rgb]{0.00,0.00,1.00}{##1}}}
\@namedef{PY@tok@nc}{\let\PY@bf=\textbf\def\PY@tc##1{\textcolor[rgb]{0.00,0.00,1.00}{##1}}}
\@namedef{PY@tok@nn}{\let\PY@bf=\textbf\def\PY@tc##1{\textcolor[rgb]{0.00,0.00,1.00}{##1}}}
\@namedef{PY@tok@ne}{\let\PY@bf=\textbf\def\PY@tc##1{\textcolor[rgb]{0.80,0.25,0.22}{##1}}}
\@namedef{PY@tok@nv}{\def\PY@tc##1{\textcolor[rgb]{0.10,0.09,0.49}{##1}}}
\@namedef{PY@tok@no}{\def\PY@tc##1{\textcolor[rgb]{0.53,0.00,0.00}{##1}}}
\@namedef{PY@tok@nl}{\def\PY@tc##1{\textcolor[rgb]{0.46,0.46,0.00}{##1}}}
\@namedef{PY@tok@ni}{\let\PY@bf=\textbf\def\PY@tc##1{\textcolor[rgb]{0.44,0.44,0.44}{##1}}}
\@namedef{PY@tok@na}{\def\PY@tc##1{\textcolor[rgb]{0.41,0.47,0.13}{##1}}}
\@namedef{PY@tok@nt}{\let\PY@bf=\textbf\def\PY@tc##1{\textcolor[rgb]{0.00,0.50,0.00}{##1}}}
\@namedef{PY@tok@nd}{\def\PY@tc##1{\textcolor[rgb]{0.67,0.13,1.00}{##1}}}
\@namedef{PY@tok@s}{\def\PY@tc##1{\textcolor[rgb]{0.73,0.13,0.13}{##1}}}
\@namedef{PY@tok@sd}{\let\PY@it=\textit\def\PY@tc##1{\textcolor[rgb]{0.73,0.13,0.13}{##1}}}
\@namedef{PY@tok@si}{\let\PY@bf=\textbf\def\PY@tc##1{\textcolor[rgb]{0.64,0.35,0.47}{##1}}}
\@namedef{PY@tok@se}{\let\PY@bf=\textbf\def\PY@tc##1{\textcolor[rgb]{0.67,0.36,0.12}{##1}}}
\@namedef{PY@tok@sr}{\def\PY@tc##1{\textcolor[rgb]{0.64,0.35,0.47}{##1}}}
\@namedef{PY@tok@ss}{\def\PY@tc##1{\textcolor[rgb]{0.10,0.09,0.49}{##1}}}
\@namedef{PY@tok@sx}{\def\PY@tc##1{\textcolor[rgb]{0.00,0.50,0.00}{##1}}}
\@namedef{PY@tok@m}{\def\PY@tc##1{\textcolor[rgb]{0.40,0.40,0.40}{##1}}}
\@namedef{PY@tok@gh}{\let\PY@bf=\textbf\def\PY@tc##1{\textcolor[rgb]{0.00,0.00,0.50}{##1}}}
\@namedef{PY@tok@gu}{\let\PY@bf=\textbf\def\PY@tc##1{\textcolor[rgb]{0.50,0.00,0.50}{##1}}}
\@namedef{PY@tok@gd}{\def\PY@tc##1{\textcolor[rgb]{0.63,0.00,0.00}{##1}}}
\@namedef{PY@tok@gi}{\def\PY@tc##1{\textcolor[rgb]{0.00,0.52,0.00}{##1}}}
\@namedef{PY@tok@gr}{\def\PY@tc##1{\textcolor[rgb]{0.89,0.00,0.00}{##1}}}
\@namedef{PY@tok@ge}{\let\PY@it=\textit}
\@namedef{PY@tok@gs}{\let\PY@bf=\textbf}
\@namedef{PY@tok@ges}{\let\PY@bf=\textbf\let\PY@it=\textit}
\@namedef{PY@tok@gp}{\let\PY@bf=\textbf\def\PY@tc##1{\textcolor[rgb]{0.00,0.00,0.50}{##1}}}
\@namedef{PY@tok@go}{\def\PY@tc##1{\textcolor[rgb]{0.44,0.44,0.44}{##1}}}
\@namedef{PY@tok@gt}{\def\PY@tc##1{\textcolor[rgb]{0.00,0.27,0.87}{##1}}}
\@namedef{PY@tok@err}{\def\PY@bc##1{{\setlength{\fboxsep}{\string -\fboxrule}\fcolorbox[rgb]{1.00,0.00,0.00}{1,1,1}{\strut ##1}}}}
\@namedef{PY@tok@kc}{\let\PY@bf=\textbf\def\PY@tc##1{\textcolor[rgb]{0.00,0.50,0.00}{##1}}}
\@namedef{PY@tok@kd}{\let\PY@bf=\textbf\def\PY@tc##1{\textcolor[rgb]{0.00,0.50,0.00}{##1}}}
\@namedef{PY@tok@kn}{\let\PY@bf=\textbf\def\PY@tc##1{\textcolor[rgb]{0.00,0.50,0.00}{##1}}}
\@namedef{PY@tok@kr}{\let\PY@bf=\textbf\def\PY@tc##1{\textcolor[rgb]{0.00,0.50,0.00}{##1}}}
\@namedef{PY@tok@bp}{\def\PY@tc##1{\textcolor[rgb]{0.00,0.50,0.00}{##1}}}
\@namedef{PY@tok@fm}{\def\PY@tc##1{\textcolor[rgb]{0.00,0.00,1.00}{##1}}}
\@namedef{PY@tok@vc}{\def\PY@tc##1{\textcolor[rgb]{0.10,0.09,0.49}{##1}}}
\@namedef{PY@tok@vg}{\def\PY@tc##1{\textcolor[rgb]{0.10,0.09,0.49}{##1}}}
\@namedef{PY@tok@vi}{\def\PY@tc##1{\textcolor[rgb]{0.10,0.09,0.49}{##1}}}
\@namedef{PY@tok@vm}{\def\PY@tc##1{\textcolor[rgb]{0.10,0.09,0.49}{##1}}}
\@namedef{PY@tok@sa}{\def\PY@tc##1{\textcolor[rgb]{0.73,0.13,0.13}{##1}}}
\@namedef{PY@tok@sb}{\def\PY@tc##1{\textcolor[rgb]{0.73,0.13,0.13}{##1}}}
\@namedef{PY@tok@sc}{\def\PY@tc##1{\textcolor[rgb]{0.73,0.13,0.13}{##1}}}
\@namedef{PY@tok@dl}{\def\PY@tc##1{\textcolor[rgb]{0.73,0.13,0.13}{##1}}}
\@namedef{PY@tok@s2}{\def\PY@tc##1{\textcolor[rgb]{0.73,0.13,0.13}{##1}}}
\@namedef{PY@tok@sh}{\def\PY@tc##1{\textcolor[rgb]{0.73,0.13,0.13}{##1}}}
\@namedef{PY@tok@s1}{\def\PY@tc##1{\textcolor[rgb]{0.73,0.13,0.13}{##1}}}
\@namedef{PY@tok@mb}{\def\PY@tc##1{\textcolor[rgb]{0.40,0.40,0.40}{##1}}}
\@namedef{PY@tok@mf}{\def\PY@tc##1{\textcolor[rgb]{0.40,0.40,0.40}{##1}}}
\@namedef{PY@tok@mh}{\def\PY@tc##1{\textcolor[rgb]{0.40,0.40,0.40}{##1}}}
\@namedef{PY@tok@mi}{\def\PY@tc##1{\textcolor[rgb]{0.40,0.40,0.40}{##1}}}
\@namedef{PY@tok@il}{\def\PY@tc##1{\textcolor[rgb]{0.40,0.40,0.40}{##1}}}
\@namedef{PY@tok@mo}{\def\PY@tc##1{\textcolor[rgb]{0.40,0.40,0.40}{##1}}}
\@namedef{PY@tok@ch}{\let\PY@it=\textit\def\PY@tc##1{\textcolor[rgb]{0.24,0.48,0.48}{##1}}}
\@namedef{PY@tok@cm}{\let\PY@it=\textit\def\PY@tc##1{\textcolor[rgb]{0.24,0.48,0.48}{##1}}}
\@namedef{PY@tok@cpf}{\let\PY@it=\textit\def\PY@tc##1{\textcolor[rgb]{0.24,0.48,0.48}{##1}}}
\@namedef{PY@tok@c1}{\let\PY@it=\textit\def\PY@tc##1{\textcolor[rgb]{0.24,0.48,0.48}{##1}}}
\@namedef{PY@tok@cs}{\let\PY@it=\textit\def\PY@tc##1{\textcolor[rgb]{0.24,0.48,0.48}{##1}}}

\def\PYZbs{\char`\\}
\def\PYZus{\char`\_}
\def\PYZob{\char`\{}
\def\PYZcb{\char`\}}
\def\PYZca{\char`\^}
\def\PYZam{\char`\&}
\def\PYZlt{\char`\<}
\def\PYZgt{\char`\>}
\def\PYZsh{\char`\#}
\def\PYZpc{\char`\%}
\def\PYZdl{\char`\$}
\def\PYZhy{\char`\-}
\def\PYZsq{\char`\'}
\def\PYZdq{\char`\"}
\def\PYZti{\char`\~}
% for compatibility with earlier versions
\def\PYZat{@}
\def\PYZlb{[}
\def\PYZrb{]}
\makeatother


    % For linebreaks inside Verbatim environment from package fancyvrb.
    \makeatletter
        \newbox\Wrappedcontinuationbox
        \newbox\Wrappedvisiblespacebox
        \newcommand*\Wrappedvisiblespace {\textcolor{red}{\textvisiblespace}}
        \newcommand*\Wrappedcontinuationsymbol {\textcolor{red}{\llap{\tiny$\m@th\hookrightarrow$}}}
        \newcommand*\Wrappedcontinuationindent {3ex }
        \newcommand*\Wrappedafterbreak {\kern\Wrappedcontinuationindent\copy\Wrappedcontinuationbox}
        % Take advantage of the already applied Pygments mark-up to insert
        % potential linebreaks for TeX processing.
        %        {, <, #, %, $, ' and ": go to next line.
        %        _, }, ^, &, >, - and ~: stay at end of broken line.
        % Use of \textquotesingle for straight quote.
        \newcommand*\Wrappedbreaksatspecials {%
            \def\PYGZus{\discretionary{\char`\_}{\Wrappedafterbreak}{\char`\_}}%
            \def\PYGZob{\discretionary{}{\Wrappedafterbreak\char`\{}{\char`\{}}%
            \def\PYGZcb{\discretionary{\char`\}}{\Wrappedafterbreak}{\char`\}}}%
            \def\PYGZca{\discretionary{\char`\^}{\Wrappedafterbreak}{\char`\^}}%
            \def\PYGZam{\discretionary{\char`\&}{\Wrappedafterbreak}{\char`\&}}%
            \def\PYGZlt{\discretionary{}{\Wrappedafterbreak\char`\<}{\char`\<}}%
            \def\PYGZgt{\discretionary{\char`\>}{\Wrappedafterbreak}{\char`\>}}%
            \def\PYGZsh{\discretionary{}{\Wrappedafterbreak\char`\#}{\char`\#}}%
            \def\PYGZpc{\discretionary{}{\Wrappedafterbreak\char`\%}{\char`\%}}%
            \def\PYGZdl{\discretionary{}{\Wrappedafterbreak\char`\$}{\char`\$}}%
            \def\PYGZhy{\discretionary{\char`\-}{\Wrappedafterbreak}{\char`\-}}%
            \def\PYGZsq{\discretionary{}{\Wrappedafterbreak\textquotesingle}{\textquotesingle}}%
            \def\PYGZdq{\discretionary{}{\Wrappedafterbreak\char`\"}{\char`\"}}%
            \def\PYGZti{\discretionary{\char`\~}{\Wrappedafterbreak}{\char`\~}}%
        }
        % Some characters . , ; ? ! / are not pygmentized.
        % This macro makes them "active" and they will insert potential linebreaks
        \newcommand*\Wrappedbreaksatpunct {%
            \lccode`\~`\.\lowercase{\def~}{\discretionary{\hbox{\char`\.}}{\Wrappedafterbreak}{\hbox{\char`\.}}}%
            \lccode`\~`\,\lowercase{\def~}{\discretionary{\hbox{\char`\,}}{\Wrappedafterbreak}{\hbox{\char`\,}}}%
            \lccode`\~`\;\lowercase{\def~}{\discretionary{\hbox{\char`\;}}{\Wrappedafterbreak}{\hbox{\char`\;}}}%
            \lccode`\~`\:\lowercase{\def~}{\discretionary{\hbox{\char`\:}}{\Wrappedafterbreak}{\hbox{\char`\:}}}%
            \lccode`\~`\?\lowercase{\def~}{\discretionary{\hbox{\char`\?}}{\Wrappedafterbreak}{\hbox{\char`\?}}}%
            \lccode`\~`\!\lowercase{\def~}{\discretionary{\hbox{\char`\!}}{\Wrappedafterbreak}{\hbox{\char`\!}}}%
            \lccode`\~`\/\lowercase{\def~}{\discretionary{\hbox{\char`\/}}{\Wrappedafterbreak}{\hbox{\char`\/}}}%
            \catcode`\.\active
            \catcode`\,\active
            \catcode`\;\active
            \catcode`\:\active
            \catcode`\?\active
            \catcode`\!\active
            \catcode`\/\active
            \lccode`\~`\~
        }
    \makeatother

    \let\OriginalVerbatim=\Verbatim
    \makeatletter
    \renewcommand{\Verbatim}[1][1]{%
        %\parskip\z@skip
        \sbox\Wrappedcontinuationbox {\Wrappedcontinuationsymbol}%
        \sbox\Wrappedvisiblespacebox {\FV@SetupFont\Wrappedvisiblespace}%
        \def\FancyVerbFormatLine ##1{\hsize\linewidth
            \vtop{\raggedright\hyphenpenalty\z@\exhyphenpenalty\z@
                \doublehyphendemerits\z@\finalhyphendemerits\z@
                \strut ##1\strut}%
        }%
        % If the linebreak is at a space, the latter will be displayed as visible
        % space at end of first line, and a continuation symbol starts next line.
        % Stretch/shrink are however usually zero for typewriter font.
        \def\FV@Space {%
            \nobreak\hskip\z@ plus\fontdimen3\font minus\fontdimen4\font
            \discretionary{\copy\Wrappedvisiblespacebox}{\Wrappedafterbreak}
            {\kern\fontdimen2\font}%
        }%

        % Allow breaks at special characters using \PYG... macros.
        \Wrappedbreaksatspecials
        % Breaks at punctuation characters . , ; ? ! and / need catcode=\active
        \OriginalVerbatim[#1,codes*=\Wrappedbreaksatpunct]%
    }
    \makeatother

    % Exact colors from NB
    \definecolor{incolor}{HTML}{303F9F}
    \definecolor{outcolor}{HTML}{D84315}
    \definecolor{cellborder}{HTML}{CFCFCF}
    \definecolor{cellbackground}{HTML}{F7F7F7}

    % prompt
    \makeatletter
    \newcommand{\boxspacing}{\kern\kvtcb@left@rule\kern\kvtcb@boxsep}
    \makeatother
    \newcommand{\prompt}[4]{
        {\ttfamily\llap{{\color{#2}[#3]:\hspace{3pt}#4}}\vspace{-\baselineskip}}
    }
    

    
    % Prevent overflowing lines due to hard-to-break entities
    \sloppy
    % Setup hyperref package
    \hypersetup{
      breaklinks=true,  % so long urls are correctly broken across lines
      colorlinks=true,
      urlcolor=urlcolor,
      linkcolor=linkcolor,
      citecolor=citecolor,
      }
    % Slightly bigger margins than the latex defaults
    
    \geometry{verbose,tmargin=1in,bmargin=1in,lmargin=1in,rmargin=1in}
    
    

\newcounter{none}
\begin{document}
    
    \maketitle
    
    

    
    \section{Atividade --- Explorando Árvores de Decisão para
Classificação}\label{atividade-explorando-uxe1rvores-de-decisuxe3o-para-classificauxe7uxe3o}

\subsection{Objetivo}\label{objetivo}

Nesta atividade, você irá resolver um problema de \textbf{classificação}
utilizando Árvores de Decisão e investigar como diferentes escolhas
influenciam o modelo.

Você deverá: - construir árvores de decisão para um conjunto de dados
real; - comparar diferentes critérios de particionamento; - observar
quais atributos são escolhidos pela árvore; - visualizar os primeiros
níveis das árvores; - investigar o efeito da profundidade sobre o
desempenho; - utilizar um mecanismo de \textbf{parada antecipada
(pre-pruning)}; - comparar desempenho e complexidade; - escolher e
justificar um modelo final.

\begin{quote}
\textbf{Importante:} todas as respostas devem ser preenchidas no próprio
notebook, nas células indicadas.
\end{quote}

    \subsection{1. Conjunto de dados}\label{conjunto-de-dados}

Será utilizado o conjunto \textbf{Bank Marketing}, disponibilizado pelo
UCI Machine Learning Repository.

https://archive.ics.uci.edu/dataset/222/bank+marketing

O problema consiste em prever se um cliente irá contratar (\texttt{yes})
ou não (\texttt{no}) um depósito a prazo após uma campanha de marketing
bancário.

O arquivo utilizado é \texttt{bank-full.csv}, com separador \texttt{;}.

O conjunto possui 45.211 instâncias e atributos numéricos e categóricos.

\textbf{Fonte:} UCI Machine Learning Repository --- Bank Marketing

O arquivo \texttt{bank-full.csv} deve estar na mesma pasta deste
notebook.

    \subsection{2. Importação das
bibliotecas}\label{importauxe7uxe3o-das-bibliotecas}

Execute a célula a seguir. O código inicial já contém as principais
bibliotecas necessárias.

    \begin{tcolorbox}[breakable, size=fbox, boxrule=1pt, pad at break*=1mm,colback=cellbackground, colframe=cellborder]
\prompt{In}{incolor}{1}{\boxspacing}
\begin{Verbatim}[commandchars=\\\{\}]
\PY{k+kn}{import}\PY{+w}{ }\PY{n+nn}{pandas}\PY{+w}{ }\PY{k}{as}\PY{+w}{ }\PY{n+nn}{pd}
\PY{k+kn}{import}\PY{+w}{ }\PY{n+nn}{numpy}\PY{+w}{ }\PY{k}{as}\PY{+w}{ }\PY{n+nn}{np}
\PY{k+kn}{import}\PY{+w}{ }\PY{n+nn}{matplotlib}\PY{n+nn}{.}\PY{n+nn}{pyplot}\PY{+w}{ }\PY{k}{as}\PY{+w}{ }\PY{n+nn}{plt}

\PY{k+kn}{from}\PY{+w}{ }\PY{n+nn}{sklearn}\PY{n+nn}{.}\PY{n+nn}{model\PYZus{}selection}\PY{+w}{ }\PY{k+kn}{import} \PY{n}{train\PYZus{}test\PYZus{}split}
\PY{k+kn}{from}\PY{+w}{ }\PY{n+nn}{sklearn}\PY{n+nn}{.}\PY{n+nn}{compose}\PY{+w}{ }\PY{k+kn}{import} \PY{n}{ColumnTransformer}
\PY{k+kn}{from}\PY{+w}{ }\PY{n+nn}{sklearn}\PY{n+nn}{.}\PY{n+nn}{preprocessing}\PY{+w}{ }\PY{k+kn}{import} \PY{n}{OneHotEncoder}
\PY{k+kn}{from}\PY{+w}{ }\PY{n+nn}{sklearn}\PY{n+nn}{.}\PY{n+nn}{tree}\PY{+w}{ }\PY{k+kn}{import} \PY{n}{DecisionTreeClassifier}\PY{p}{,} \PY{n}{plot\PYZus{}tree}
\PY{k+kn}{from}\PY{+w}{ }\PY{n+nn}{sklearn}\PY{n+nn}{.}\PY{n+nn}{metrics}\PY{+w}{ }\PY{k+kn}{import} \PY{p}{(}
    \PY{n}{accuracy\PYZus{}score}\PY{p}{,} \PY{n}{precision\PYZus{}score}\PY{p}{,} \PY{n}{recall\PYZus{}score}\PY{p}{,} \PY{n}{f1\PYZus{}score}\PY{p}{,}
    \PY{n}{classification\PYZus{}report}\PY{p}{,} \PY{n}{ConfusionMatrixDisplay}
\PY{p}{)}
\end{Verbatim}
\end{tcolorbox}

    \subsection{3. Carregamento dos dados}\label{carregamento-dos-dados}

Execute o código abaixo e observe as primeiras linhas.

    \begin{tcolorbox}[breakable, size=fbox, boxrule=1pt, pad at break*=1mm,colback=cellbackground, colframe=cellborder]
\prompt{In}{incolor}{2}{\boxspacing}
\begin{Verbatim}[commandchars=\\\{\}]
\PY{n}{df} \PY{o}{=} \PY{n}{pd}\PY{o}{.}\PY{n}{read\PYZus{}csv}\PY{p}{(}\PY{l+s+s2}{\PYZdq{}}\PY{l+s+s2}{bank\PYZhy{}full.csv}\PY{l+s+s2}{\PYZdq{}}\PY{p}{,} \PY{n}{sep}\PY{o}{=}\PY{l+s+s2}{\PYZdq{}}\PY{l+s+s2}{;}\PY{l+s+s2}{\PYZdq{}}\PY{p}{)}
\PY{n}{df}\PY{o}{.}\PY{n}{head}\PY{p}{(}\PY{p}{)}
\end{Verbatim}
\end{tcolorbox}

            \begin{tcolorbox}[breakable, size=fbox, boxrule=.5pt, pad at break*=1mm, opacityfill=0]
\prompt{Out}{outcolor}{2}{\boxspacing}
\begin{Verbatim}[commandchars=\\\{\}]
   age           job  marital  education default  balance housing loan  \textbackslash{}
0   58    management  married   tertiary      no     2143     yes   no
1   44    technician   single  secondary      no       29     yes   no
2   33  entrepreneur  married  secondary      no        2     yes  yes
3   47   blue-collar  married    unknown      no     1506     yes   no
4   33       unknown   single    unknown      no        1      no   no

   contact  day month  duration  campaign  pdays  previous poutcome   y
0  unknown    5   may       261         1     -1         0  unknown  no
1  unknown    5   may       151         1     -1         0  unknown  no
2  unknown    5   may        76         1     -1         0  unknown  no
3  unknown    5   may        92         1     -1         0  unknown  no
4  unknown    5   may       198         1     -1         0  unknown  no
\end{Verbatim}
\end{tcolorbox}
        
    \subsubsection{Questão 1 --- Conhecendo os
dados}\label{questuxe3o-1-conhecendo-os-dados}

Investigue o conjunto utilizando \texttt{shape}, \texttt{info},
\texttt{describe} e outros comandos que considerar adequados.

Responda: 1. Quantas instâncias existem? 2. Quantos atributos existem?
3. Qual é a variável-alvo? 4. Quais atributos são numéricos? 5. Quais
atributos são categóricos?

\textbf{Resposta:}

    \begin{tcolorbox}[breakable, size=fbox, boxrule=1pt, pad at break*=1mm,colback=cellbackground, colframe=cellborder]
\prompt{In}{incolor}{3}{\boxspacing}
\begin{Verbatim}[commandchars=\\\{\}]
\PY{n+nb}{print}\PY{p}{(}\PY{l+s+s1}{\PYZsq{}}\PY{l+s+s1}{Shape:}\PY{l+s+s1}{\PYZsq{}}\PY{p}{,} \PY{n}{df}\PY{o}{.}\PY{n}{shape}\PY{p}{)}
\PY{n}{df}\PY{o}{.}\PY{n}{info}\PY{p}{(}\PY{p}{)}
\PY{n}{df}\PY{o}{.}\PY{n}{describe}\PY{p}{(}\PY{n}{include}\PY{o}{=}\PY{l+s+s1}{\PYZsq{}}\PY{l+s+s1}{all}\PY{l+s+s1}{\PYZsq{}}\PY{p}{)}
\end{Verbatim}
\end{tcolorbox}

    \begin{Verbatim}[commandchars=\\\{\}]
Shape: (45211, 17)
<class 'pandas.DataFrame'>
RangeIndex: 45211 entries, 0 to 45210
Data columns (total 17 columns):
 \#   Column     Non-Null Count  Dtype
---  ------     --------------  -----
 0   age        45211 non-null  int64
 1   job        45211 non-null  str
 2   marital    45211 non-null  str
 3   education  45211 non-null  str
 4   default    45211 non-null  str
 5   balance    45211 non-null  int64
 6   housing    45211 non-null  str
 7   loan       45211 non-null  str
 8   contact    45211 non-null  str
 9   day        45211 non-null  int64
 10  month      45211 non-null  str
 11  duration   45211 non-null  int64
 12  campaign   45211 non-null  int64
 13  pdays      45211 non-null  int64
 14  previous   45211 non-null  int64
 15  poutcome   45211 non-null  str
 16  y          45211 non-null  str
dtypes: int64(7), str(10)
memory usage: 5.9 MB
    \end{Verbatim}

            \begin{tcolorbox}[breakable, size=fbox, boxrule=.5pt, pad at break*=1mm, opacityfill=0]
\prompt{Out}{outcolor}{3}{\boxspacing}
\begin{Verbatim}[commandchars=\\\{\}]
                 age          job  marital  education default        balance  \textbackslash{}
count   45211.000000        45211    45211      45211   45211   45211.000000
unique           NaN           12        3          4       2            NaN
top              NaN  blue-collar  married  secondary      no            NaN
freq             NaN         9732    27214      23202   44396            NaN
mean       40.936210          NaN      NaN        NaN     NaN    1362.272058
std        10.618762          NaN      NaN        NaN     NaN    3044.765829
min        18.000000          NaN      NaN        NaN     NaN   -8019.000000
25\%        33.000000          NaN      NaN        NaN     NaN      72.000000
50\%        39.000000          NaN      NaN        NaN     NaN     448.000000
75\%        48.000000          NaN      NaN        NaN     NaN    1428.000000
max        95.000000          NaN      NaN        NaN     NaN  102127.000000

       housing   loan   contact           day  month      duration  \textbackslash{}
count    45211  45211     45211  45211.000000  45211  45211.000000
unique       2      2         3           NaN     12           NaN
top        yes     no  cellular           NaN    may           NaN
freq     25130  37967     29285           NaN  13766           NaN
mean       NaN    NaN       NaN     15.806419    NaN    258.163080
std        NaN    NaN       NaN      8.322476    NaN    257.527812
min        NaN    NaN       NaN      1.000000    NaN      0.000000
25\%        NaN    NaN       NaN      8.000000    NaN    103.000000
50\%        NaN    NaN       NaN     16.000000    NaN    180.000000
75\%        NaN    NaN       NaN     21.000000    NaN    319.000000
max        NaN    NaN       NaN     31.000000    NaN   4918.000000

            campaign         pdays      previous poutcome      y
count   45211.000000  45211.000000  45211.000000    45211  45211
unique           NaN           NaN           NaN        4      2
top              NaN           NaN           NaN  unknown     no
freq             NaN           NaN           NaN    36959  39922
mean        2.763841     40.197828      0.580323      NaN    NaN
std         3.098021    100.128746      2.303441      NaN    NaN
min         1.000000     -1.000000      0.000000      NaN    NaN
25\%         1.000000     -1.000000      0.000000      NaN    NaN
50\%         2.000000     -1.000000      0.000000      NaN    NaN
75\%         3.000000     -1.000000      0.000000      NaN    NaN
max        63.000000    871.000000    275.000000      NaN    NaN
\end{Verbatim}
\end{tcolorbox}
        
    Resposta:

\begin{enumerate}
\def\labelenumi{\arabic{enumi}.}
\tightlist
\item
  45.211 instâncias
\item
  17 atributos no total (16 previsores + 1 variável-alvo)
\item
  y (indica se o cliente contratou o depósito a prazo: `yes' ou `no')
\item
  age, balance, day, duration, campaign, pdays, previous (7 atributos)
\item
  job, marital, education, default, housing, loan, contact, month,
  poutcome (9 atributos previsores) e y (alvo)
\end{enumerate}

    \subsection{4. Distribuição das
classes}\label{distribuiuxe7uxe3o-das-classes}

Investigue a distribuição da variável-alvo.

    \begin{tcolorbox}[breakable, size=fbox, boxrule=1pt, pad at break*=1mm,colback=cellbackground, colframe=cellborder]
\prompt{In}{incolor}{4}{\boxspacing}
\begin{Verbatim}[commandchars=\\\{\}]
\PY{n}{distY}\PY{o}{=} \PY{n}{df}\PY{p}{[}\PY{l+s+s2}{\PYZdq{}}\PY{l+s+s2}{y}\PY{l+s+s2}{\PYZdq{}}\PY{p}{]}\PY{o}{.}\PY{n}{value\PYZus{}counts}\PY{p}{(}\PY{p}{)}
\end{Verbatim}
\end{tcolorbox}

    \begin{tcolorbox}[breakable, size=fbox, boxrule=1pt, pad at break*=1mm,colback=cellbackground, colframe=cellborder]
\prompt{In}{incolor}{5}{\boxspacing}
\begin{Verbatim}[commandchars=\\\{\}]
\PY{n}{distY}\PY{p}{[}\PY{l+s+s1}{\PYZsq{}}\PY{l+s+s1}{yes}\PY{l+s+s1}{\PYZsq{}}\PY{p}{]}\PY{o}{/}\PY{n}{distY}\PY{p}{[}\PY{l+s+s1}{\PYZsq{}}\PY{l+s+s1}{no}\PY{l+s+s1}{\PYZsq{}}\PY{p}{]}
\end{Verbatim}
\end{tcolorbox}

            \begin{tcolorbox}[breakable, size=fbox, boxrule=.5pt, pad at break*=1mm, opacityfill=0]
\prompt{Out}{outcolor}{5}{\boxspacing}
\begin{Verbatim}[commandchars=\\\{\}]
np.float64(0.1324833425179099)
\end{Verbatim}
\end{tcolorbox}
        
    \begin{tcolorbox}[breakable, size=fbox, boxrule=1pt, pad at break*=1mm,colback=cellbackground, colframe=cellborder]
\prompt{In}{incolor}{6}{\boxspacing}
\begin{Verbatim}[commandchars=\\\{\}]
\PY{n}{df}\PY{p}{[}\PY{l+s+s2}{\PYZdq{}}\PY{l+s+s2}{y}\PY{l+s+s2}{\PYZdq{}}\PY{p}{]}\PY{o}{.}\PY{n}{value\PYZus{}counts}\PY{p}{(}\PY{n}{normalize}\PY{o}{=}\PY{k+kc}{True}\PY{p}{)}
\end{Verbatim}
\end{tcolorbox}

            \begin{tcolorbox}[breakable, size=fbox, boxrule=.5pt, pad at break*=1mm, opacityfill=0]
\prompt{Out}{outcolor}{6}{\boxspacing}
\begin{Verbatim}[commandchars=\\\{\}]
y
no     0.883015
yes    0.116985
Name: proportion, dtype: float64
\end{Verbatim}
\end{tcolorbox}
        
    \subsubsection{Questão 2 --- Distribuição das
classes}\label{questuxe3o-2-distribuiuxe7uxe3o-das-classes}

\begin{enumerate}
\def\labelenumi{\arabic{enumi}.}
\tightlist
\item
  As classes estão balanceadas?
\item
  Qual é a classe majoritária?
\item
  Por que o desbalanceamento deve ser considerado na avaliação de um
  classificador?
\end{enumerate}

\textbf{Resposta:}

    Resposta:

\begin{enumerate}
\def\labelenumi{\arabic{enumi}.}
\tightlist
\item
  Não, as classes são bem desbalanceadas.
\item
  A classe majoritária é `no' com 88,30\% (39.922 instâncias), enquanto
  `yes' tem apenas 11,70\% (5.289 instâncias).
\item
  Porque se um classificador simplesmente chutar sempre a classe
  majoritária (`no'), ele já atinge 88,3\% de acurácia sem identificar
  nenhum cliente positivo (`yes'). Por isso precisamos analisar métricas
  como precisão, recall e F1-score.
\end{enumerate}

    \subsection{5. Preparação dos dados}\label{preparauxe7uxe3o-dos-dados}

Separe os atributos da variável-alvo. Os atributos categóricos serão
transformados em uma representação numérica por One-Hot Encoding.

    \begin{tcolorbox}[breakable, size=fbox, boxrule=1pt, pad at break*=1mm,colback=cellbackground, colframe=cellborder]
\prompt{In}{incolor}{7}{\boxspacing}
\begin{Verbatim}[commandchars=\\\{\}]
\PY{n}{X} \PY{o}{=} \PY{n}{df}\PY{o}{.}\PY{n}{drop}\PY{p}{(}\PY{l+s+s2}{\PYZdq{}}\PY{l+s+s2}{y}\PY{l+s+s2}{\PYZdq{}}\PY{p}{,} \PY{n}{axis}\PY{o}{=}\PY{l+m+mi}{1}\PY{p}{)}
\PY{n}{y} \PY{o}{=} \PY{n}{df}\PY{p}{[}\PY{l+s+s2}{\PYZdq{}}\PY{l+s+s2}{y}\PY{l+s+s2}{\PYZdq{}}\PY{p}{]}\PY{o}{.}\PY{n}{map}\PY{p}{(}\PY{p}{\PYZob{}}\PY{l+s+s2}{\PYZdq{}}\PY{l+s+s2}{no}\PY{l+s+s2}{\PYZdq{}}\PY{p}{:} \PY{l+m+mi}{0}\PY{p}{,} \PY{l+s+s2}{\PYZdq{}}\PY{l+s+s2}{yes}\PY{l+s+s2}{\PYZdq{}}\PY{p}{:} \PY{l+m+mi}{1}\PY{p}{\PYZcb{}}\PY{p}{)}

\PY{n}{categorical\PYZus{}features} \PY{o}{=} \PY{n}{X}\PY{o}{.}\PY{n}{select\PYZus{}dtypes}\PY{p}{(}\PY{n}{include}\PY{o}{=}\PY{p}{[}\PY{l+s+s2}{\PYZdq{}}\PY{l+s+s2}{object}\PY{l+s+s2}{\PYZdq{}}\PY{p}{]}\PY{p}{)}\PY{o}{.}\PY{n}{columns}
\PY{n}{numeric\PYZus{}features} \PY{o}{=} \PY{n}{X}\PY{o}{.}\PY{n}{select\PYZus{}dtypes}\PY{p}{(}\PY{n}{exclude}\PY{o}{=}\PY{p}{[}\PY{l+s+s2}{\PYZdq{}}\PY{l+s+s2}{object}\PY{l+s+s2}{\PYZdq{}}\PY{p}{]}\PY{p}{)}\PY{o}{.}\PY{n}{columns}

\PY{n+nb}{print}\PY{p}{(}\PY{l+s+s2}{\PYZdq{}}\PY{l+s+s2}{Categóricos:}\PY{l+s+s2}{\PYZdq{}}\PY{p}{,} \PY{n+nb}{list}\PY{p}{(}\PY{n}{categorical\PYZus{}features}\PY{p}{)}\PY{p}{)}
\PY{n+nb}{print}\PY{p}{(}\PY{l+s+s2}{\PYZdq{}}\PY{l+s+s2}{Numéricos:}\PY{l+s+s2}{\PYZdq{}}\PY{p}{,} \PY{n+nb}{list}\PY{p}{(}\PY{n}{numeric\PYZus{}features}\PY{p}{)}\PY{p}{)}
\end{Verbatim}
\end{tcolorbox}

    \begin{Verbatim}[commandchars=\\\{\}]
Categóricos: ['job', 'marital', 'education', 'default', 'housing', 'loan',
'contact', 'month', 'poutcome']
Numéricos: ['age', 'balance', 'day', 'duration', 'campaign', 'pdays',
'previous']
    \end{Verbatim}

    \begin{Verbatim}[commandchars=\\\{\}]
/tmp/ipykernel\_467984/2409206470.py:4: Pandas4Warning: For backward
compatibility, 'str' dtypes are included by select\_dtypes when 'object' dtype is
specified. This behavior is deprecated and will be removed in a future version.
Explicitly pass 'str' to `include` to select them, or to `exclude` to remove
them and silence this warning.
See https://pandas.pydata.org/docs/user\_guide/migration-3-strings.html\#string-
migration-select-dtypes for details on how to write code that works with pandas
2 and 3.
  categorical\_features = X.select\_dtypes(include=["object"]).columns
    \end{Verbatim}

    \begin{tcolorbox}[breakable, size=fbox, boxrule=1pt, pad at break*=1mm,colback=cellbackground, colframe=cellborder]
\prompt{In}{incolor}{8}{\boxspacing}
\begin{Verbatim}[commandchars=\\\{\}]
\PY{n}{preprocessor} \PY{o}{=} \PY{n}{ColumnTransformer}\PY{p}{(}
    \PY{n}{transformers}\PY{o}{=}\PY{p}{[}
        \PY{p}{(}\PY{l+s+s2}{\PYZdq{}}\PY{l+s+s2}{cat}\PY{l+s+s2}{\PYZdq{}}\PY{p}{,} \PY{n}{OneHotEncoder}\PY{p}{(}\PY{n}{handle\PYZus{}unknown}\PY{o}{=}\PY{l+s+s2}{\PYZdq{}}\PY{l+s+s2}{ignore}\PY{l+s+s2}{\PYZdq{}}\PY{p}{)}\PY{p}{,} \PY{n}{categorical\PYZus{}features}\PY{p}{)}
    \PY{p}{]}\PY{p}{,}
    \PY{n}{remainder}\PY{o}{=}\PY{l+s+s2}{\PYZdq{}}\PY{l+s+s2}{passthrough}\PY{l+s+s2}{\PYZdq{}}
\PY{p}{)}

\PY{n}{X\PYZus{}encoded} \PY{o}{=} \PY{n}{preprocessor}\PY{o}{.}\PY{n}{fit\PYZus{}transform}\PY{p}{(}\PY{n}{X}\PY{p}{)}

\PY{n+nb}{print}\PY{p}{(}\PY{l+s+s2}{\PYZdq{}}\PY{l+s+s2}{Dimensão original:}\PY{l+s+s2}{\PYZdq{}}\PY{p}{,} \PY{n}{X}\PY{o}{.}\PY{n}{shape}\PY{p}{)}
\PY{n+nb}{print}\PY{p}{(}\PY{l+s+s2}{\PYZdq{}}\PY{l+s+s2}{Dimensão após codificação:}\PY{l+s+s2}{\PYZdq{}}\PY{p}{,} \PY{n}{X\PYZus{}encoded}\PY{o}{.}\PY{n}{shape}\PY{p}{)}
\PY{n+nb}{print}\PY{p}{(}\PY{n}{X}\PY{o}{.}\PY{n}{head}\PY{p}{(}\PY{p}{)}\PY{p}{)}
\end{Verbatim}
\end{tcolorbox}

    \begin{Verbatim}[commandchars=\\\{\}]
Dimensão original: (45211, 16)
Dimensão após codificação: (45211, 51)
   age           job  marital  education default  balance housing loan  \textbackslash{}
0   58    management  married   tertiary      no     2143     yes   no
1   44    technician   single  secondary      no       29     yes   no
2   33  entrepreneur  married  secondary      no        2     yes  yes
3   47   blue-collar  married    unknown      no     1506     yes   no
4   33       unknown   single    unknown      no        1      no   no

   contact  day month  duration  campaign  pdays  previous poutcome
0  unknown    5   may       261         1     -1         0  unknown
1  unknown    5   may       151         1     -1         0  unknown
2  unknown    5   may        76         1     -1         0  unknown
3  unknown    5   may        92         1     -1         0  unknown
4  unknown    5   may       198         1     -1         0  unknown
    \end{Verbatim}

    \subsubsection{Questão 3 --- Explique o código acima faz. Como funciona
o
OneHotEncoder?}\label{questuxe3o-3-explique-o-cuxf3digo-acima-faz.-como-funciona-o-onehotencoder}

Resposta:

\begin{itemize}
\tightlist
\item
  O que o código faz: Separa as features em X e a variável-alvo em y
  (mapeando `no'-\textgreater0 e `yes'-\textgreater1). Identifica quais
  colunas são categóricas e numéricas e usa o ColumnTransformer para
  aplicar o OneHotEncoder nas colunas categóricas e manter as numéricas
  inalteradas (remainder=`passthrough'), gerando uma matriz final com 51
  atributos numéricos.
\item
  OneHotEncoderc ria uma colunbinária (0 ou 1) para cada categoria única
  de um atributo. A categoria correspondente recebe 1 e as outras 0.
  Isso evita que o modelo atribua uma ordem de grandeza artificial entre
  categorias nominais.
\end{itemize}

    \subsection{6. Separação entre treinamento e
teste}\label{separauxe7uxe3o-entre-treinamento-e-teste}

Utilize 70\% dos dados para treinamento e 30\% para teste.

    \begin{tcolorbox}[breakable, size=fbox, boxrule=1pt, pad at break*=1mm,colback=cellbackground, colframe=cellborder]
\prompt{In}{incolor}{9}{\boxspacing}
\begin{Verbatim}[commandchars=\\\{\}]
\PY{n}{X\PYZus{}train}\PY{p}{,} \PY{n}{X\PYZus{}test}\PY{p}{,} \PY{n}{y\PYZus{}train}\PY{p}{,} \PY{n}{y\PYZus{}test} \PY{o}{=} \PY{n}{train\PYZus{}test\PYZus{}split}\PY{p}{(}
    \PY{n}{X\PYZus{}encoded}\PY{p}{,} \PY{n}{y}\PY{p}{,}
    \PY{n}{test\PYZus{}size}\PY{o}{=}\PY{l+m+mf}{0.30}\PY{p}{,}
    \PY{n}{random\PYZus{}state}\PY{o}{=}\PY{l+m+mi}{42}\PY{p}{,}
    \PY{n}{stratify}\PY{o}{=}\PY{n}{y}
\PY{p}{)}

\PY{n+nb}{print}\PY{p}{(}\PY{l+s+s2}{\PYZdq{}}\PY{l+s+s2}{Treinamento:}\PY{l+s+s2}{\PYZdq{}}\PY{p}{,} \PY{n}{X\PYZus{}train}\PY{o}{.}\PY{n}{shape}\PY{p}{)}
\PY{n+nb}{print}\PY{p}{(}\PY{l+s+s2}{\PYZdq{}}\PY{l+s+s2}{Teste:}\PY{l+s+s2}{\PYZdq{}}\PY{p}{,} \PY{n}{X\PYZus{}test}\PY{o}{.}\PY{n}{shape}\PY{p}{)}
\end{Verbatim}
\end{tcolorbox}

    \begin{Verbatim}[commandchars=\\\{\}]
Treinamento: (31647, 51)
Teste: (13564, 51)
    \end{Verbatim}

    \subsubsection{Questão 4 --- Treinamento e
teste}\label{questuxe3o-4-treinamento-e-teste}

Explique por que o conjunto de teste não deve ser utilizado para ajustar
o modelo.

Em sua resposta, explique a relação entre treinamento, teste,
generalização e overfitting.

\textbf{Resposta:}

    Resposta:

\begin{itemize}
\tightlist
\item
  Por que não usar o teste no ajuste: O conjunto de teste serve para
  simular dados novos que o modelo nunca viu. Se usarmos o teste para
  ajustar o modelo, ocorre vazamento de dados (data leakage) e a
  avaliação fica enviesada e irreal.
\item
  Relação entre os conceitos:

  \begin{itemize}
  \tightlist
  \item
    Treinamento: onde o modelo ajusta seus parâmetros e aprende as
    regras da árvore.
  \item
    Generalização: capacidade do modelo de acertar previsões em novos
    dados reais.
  \item
    Overfitting: quando o modelo decora os dados de treino (acurácia
    \textasciitilde100\% no treino), mas tem desempenho ruim no teste.
  \item
    Teste: serve justamente para medir a capacidade real de
    generalização e detectar se houve overfitting.
  \end{itemize}
\end{itemize}

    \section{Parte I --- Critérios de
particionamento}\label{parte-i-crituxe9rios-de-particionamento}

\subsection{7. Árvore de Decisão com
Gini}\label{uxe1rvore-de-decisuxe3o-com-gini}

Construa uma árvore utilizando o critério gini.

Inicialmente, não limite a profundidade. O objetivo é observar uma
árvore que cresça livremente.

    \begin{tcolorbox}[breakable, size=fbox, boxrule=1pt, pad at break*=1mm,colback=cellbackground, colframe=cellborder]
\prompt{In}{incolor}{10}{\boxspacing}
\begin{Verbatim}[commandchars=\\\{\}]
\PY{n}{tree\PYZus{}gini} \PY{o}{=} \PY{n}{DecisionTreeClassifier}\PY{p}{(}
    \PY{n}{criterion}\PY{o}{=}\PY{l+s+s2}{\PYZdq{}}\PY{l+s+s2}{gini}\PY{l+s+s2}{\PYZdq{}}\PY{p}{,}
    \PY{n}{random\PYZus{}state}\PY{o}{=}\PY{l+m+mi}{42}
\PY{p}{)}

\PY{n}{tree\PYZus{}gini}\PY{o}{.}\PY{n}{fit}\PY{p}{(}\PY{n}{X\PYZus{}train}\PY{p}{,} \PY{n}{y\PYZus{}train}\PY{p}{)}
\end{Verbatim}
\end{tcolorbox}

            \begin{tcolorbox}[breakable, size=fbox, boxrule=.5pt, pad at break*=1mm, opacityfill=0]
\prompt{Out}{outcolor}{10}{\boxspacing}
\begin{Verbatim}[commandchars=\\\{\}]
DecisionTreeClassifier(random\_state=42)
\end{Verbatim}
\end{tcolorbox}
        
    \begin{tcolorbox}[breakable, size=fbox, boxrule=1pt, pad at break*=1mm,colback=cellbackground, colframe=cellborder]
\prompt{In}{incolor}{11}{\boxspacing}
\begin{Verbatim}[commandchars=\\\{\}]
\PY{n}{y\PYZus{}pred\PYZus{}gini} \PY{o}{=} \PY{n}{tree\PYZus{}gini}\PY{o}{.}\PY{n}{predict}\PY{p}{(}\PY{n}{X\PYZus{}test}\PY{p}{)}

\PY{n+nb}{print}\PY{p}{(}\PY{l+s+s2}{\PYZdq{}}\PY{l+s+s2}{Acurácia:}\PY{l+s+s2}{\PYZdq{}}\PY{p}{,} \PY{n}{accuracy\PYZus{}score}\PY{p}{(}\PY{n}{y\PYZus{}test}\PY{p}{,} \PY{n}{y\PYZus{}pred\PYZus{}gini}\PY{p}{)}\PY{p}{)}
\PY{n+nb}{print}\PY{p}{(}\PY{l+s+s2}{\PYZdq{}}\PY{l+s+s2}{Precisão:}\PY{l+s+s2}{\PYZdq{}}\PY{p}{,} \PY{n}{precision\PYZus{}score}\PY{p}{(}\PY{n}{y\PYZus{}test}\PY{p}{,} \PY{n}{y\PYZus{}pred\PYZus{}gini}\PY{p}{)}\PY{p}{)}
\PY{n+nb}{print}\PY{p}{(}\PY{l+s+s2}{\PYZdq{}}\PY{l+s+s2}{Recall:}\PY{l+s+s2}{\PYZdq{}}\PY{p}{,} \PY{n}{recall\PYZus{}score}\PY{p}{(}\PY{n}{y\PYZus{}test}\PY{p}{,} \PY{n}{y\PYZus{}pred\PYZus{}gini}\PY{p}{)}\PY{p}{)}
\PY{n+nb}{print}\PY{p}{(}\PY{l+s+s2}{\PYZdq{}}\PY{l+s+s2}{F1:}\PY{l+s+s2}{\PYZdq{}}\PY{p}{,} \PY{n}{f1\PYZus{}score}\PY{p}{(}\PY{n}{y\PYZus{}test}\PY{p}{,} \PY{n}{y\PYZus{}pred\PYZus{}gini}\PY{p}{)}\PY{p}{)}
\PY{n+nb}{print}\PY{p}{(}\PY{l+s+s2}{\PYZdq{}}\PY{l+s+s2}{Profundidade:}\PY{l+s+s2}{\PYZdq{}}\PY{p}{,} \PY{n}{tree\PYZus{}gini}\PY{o}{.}\PY{n}{get\PYZus{}depth}\PY{p}{(}\PY{p}{)}\PY{p}{)}
\PY{n+nb}{print}\PY{p}{(}\PY{l+s+s2}{\PYZdq{}}\PY{l+s+s2}{Número de folhas:}\PY{l+s+s2}{\PYZdq{}}\PY{p}{,} \PY{n}{tree\PYZus{}gini}\PY{o}{.}\PY{n}{get\PYZus{}n\PYZus{}leaves}\PY{p}{(}\PY{p}{)}\PY{p}{)}
\end{Verbatim}
\end{tcolorbox}

    \begin{Verbatim}[commandchars=\\\{\}]
Acurácia: 0.8731200235918608
Precisão: 0.45934466019417475
Recall: 0.47700063011972277
F1: 0.4680061823802164
Profundidade: 36
Número de folhas: 2562
    \end{Verbatim}

    \subsection{8. Visualização da
árvore}\label{visualizauxe7uxe3o-da-uxe1rvore}

    \begin{tcolorbox}[breakable, size=fbox, boxrule=1pt, pad at break*=1mm,colback=cellbackground, colframe=cellborder]
\prompt{In}{incolor}{12}{\boxspacing}
\begin{Verbatim}[commandchars=\\\{\}]
\PY{n}{feature\PYZus{}names} \PY{o}{=} \PY{n}{preprocessor}\PY{o}{.}\PY{n}{get\PYZus{}feature\PYZus{}names\PYZus{}out}\PY{p}{(}\PY{p}{)}

\PY{n}{plt}\PY{o}{.}\PY{n}{figure}\PY{p}{(}\PY{n}{figsize}\PY{o}{=}\PY{p}{(}\PY{l+m+mi}{24}\PY{p}{,} \PY{l+m+mi}{12}\PY{p}{)}\PY{p}{)}
\PY{n}{plot\PYZus{}tree}\PY{p}{(}
    \PY{n}{tree\PYZus{}gini}\PY{p}{,}
    \PY{n}{max\PYZus{}depth}\PY{o}{=}\PY{l+m+mi}{10}\PY{p}{,}
    \PY{n}{feature\PYZus{}names}\PY{o}{=}\PY{n}{feature\PYZus{}names}\PY{p}{,}
    \PY{n}{class\PYZus{}names}\PY{o}{=}\PY{p}{[}\PY{l+s+s2}{\PYZdq{}}\PY{l+s+s2}{no}\PY{l+s+s2}{\PYZdq{}}\PY{p}{,} \PY{l+s+s2}{\PYZdq{}}\PY{l+s+s2}{yes}\PY{l+s+s2}{\PYZdq{}}\PY{p}{]}\PY{p}{,}
    \PY{n}{filled}\PY{o}{=}\PY{k+kc}{True}\PY{p}{,}
    \PY{n}{fontsize}\PY{o}{=}\PY{l+m+mi}{8}
\PY{p}{)}
\PY{n}{plt}\PY{o}{.}\PY{n}{show}\PY{p}{(}\PY{p}{)}
\end{Verbatim}
\end{tcolorbox}

    \begin{center}
    \adjustimage{max size={0.9\linewidth}{0.9\paperheight}}{output_27_0.png}
    \end{center}
    { \hspace*{\fill} \\}
    
    \subsubsection{Questão 5 --- Observando os
particionamentos}\label{questuxe3o-5-observando-os-particionamentos}

Observe os primeiros níveis da árvore.

\begin{enumerate}
\def\labelenumi{\arabic{enumi}.}
\tightlist
\item
  Qual atributo foi escolhido na raiz?
\item
  Qual condição de particionamento foi utilizada?
\item
  Quais atributos aparecem nos níveis seguintes?
\item
  O atributo escolhido na raiz é necessariamente o atributo mais
  importante de todo o modelo? Explique.
\item
  O que representam gini, samples, value e class em um nó?
\end{enumerate}

Resposta:

    \textbf{Resposta:}

\begin{enumerate}
\def\labelenumi{\arabic{enumi}.}
\tightlist
\item
  remainder\_\_duration (duração do último contato em segundos).
\item
  duration \textless= 518.5
\item
  Nos níveis seguintes aparecem cat\_\_poutcome\_success, duration
  \textless= 827.5, age, day, month\_oct, month\_mar, pdays,
  housing\_yes, etc.
\item
  Não necessariamente. A árvore escolhe na raiz o melhor atributo de
  forma local/imediata (algoritmo guloso). Mas atributos que aparecem
  várias vezes ao longo da árvore podem somar mais importância no modelo
  geral.
\item
  \begin{itemize}
  \tightlist
  \item
    gini: nível de impureza do nó (0 significa nó puro);
  \item
    samples: número de exemplos de treino que chegaram ao nó;
  \item
    value: contagem de amostras de cada classe no nó ({[}no, yes{]});
  \item
    class: classe majoritária naquele nó.
  \end{itemize}
\end{enumerate}

    \subsection{9. Árvore de Decisão com
Entropy}\label{uxe1rvore-de-decisuxe3o-com-entropy}

Agora construa outra árvore, utilizando entropy como critério.

    \begin{tcolorbox}[breakable, size=fbox, boxrule=1pt, pad at break*=1mm,colback=cellbackground, colframe=cellborder]
\prompt{In}{incolor}{13}{\boxspacing}
\begin{Verbatim}[commandchars=\\\{\}]
\PY{n}{tree\PYZus{}entropy} \PY{o}{=} \PY{n}{DecisionTreeClassifier}\PY{p}{(}
    \PY{n}{criterion}\PY{o}{=}\PY{l+s+s2}{\PYZdq{}}\PY{l+s+s2}{entropy}\PY{l+s+s2}{\PYZdq{}}\PY{p}{,}
    \PY{n}{random\PYZus{}state}\PY{o}{=}\PY{l+m+mi}{42}
\PY{p}{)}

\PY{n}{tree\PYZus{}entropy}\PY{o}{.}\PY{n}{fit}\PY{p}{(}\PY{n}{X\PYZus{}train}\PY{p}{,} \PY{n}{y\PYZus{}train}\PY{p}{)}

\PY{n}{y\PYZus{}pred\PYZus{}entropy} \PY{o}{=} \PY{n}{tree\PYZus{}entropy}\PY{o}{.}\PY{n}{predict}\PY{p}{(}\PY{n}{X\PYZus{}test}\PY{p}{)}

\PY{n+nb}{print}\PY{p}{(}\PY{l+s+s2}{\PYZdq{}}\PY{l+s+s2}{Acurácia:}\PY{l+s+s2}{\PYZdq{}}\PY{p}{,} \PY{n}{accuracy\PYZus{}score}\PY{p}{(}\PY{n}{y\PYZus{}test}\PY{p}{,} \PY{n}{y\PYZus{}pred\PYZus{}entropy}\PY{p}{)}\PY{p}{)}
\PY{n+nb}{print}\PY{p}{(}\PY{l+s+s2}{\PYZdq{}}\PY{l+s+s2}{Precisão:}\PY{l+s+s2}{\PYZdq{}}\PY{p}{,} \PY{n}{precision\PYZus{}score}\PY{p}{(}\PY{n}{y\PYZus{}test}\PY{p}{,} \PY{n}{y\PYZus{}pred\PYZus{}entropy}\PY{p}{)}\PY{p}{)}
\PY{n+nb}{print}\PY{p}{(}\PY{l+s+s2}{\PYZdq{}}\PY{l+s+s2}{Recall:}\PY{l+s+s2}{\PYZdq{}}\PY{p}{,} \PY{n}{recall\PYZus{}score}\PY{p}{(}\PY{n}{y\PYZus{}test}\PY{p}{,} \PY{n}{y\PYZus{}pred\PYZus{}entropy}\PY{p}{)}\PY{p}{)}
\PY{n+nb}{print}\PY{p}{(}\PY{l+s+s2}{\PYZdq{}}\PY{l+s+s2}{F1:}\PY{l+s+s2}{\PYZdq{}}\PY{p}{,} \PY{n}{f1\PYZus{}score}\PY{p}{(}\PY{n}{y\PYZus{}test}\PY{p}{,} \PY{n}{y\PYZus{}pred\PYZus{}entropy}\PY{p}{)}\PY{p}{)}
\PY{n+nb}{print}\PY{p}{(}\PY{l+s+s2}{\PYZdq{}}\PY{l+s+s2}{Profundidade:}\PY{l+s+s2}{\PYZdq{}}\PY{p}{,} \PY{n}{tree\PYZus{}entropy}\PY{o}{.}\PY{n}{get\PYZus{}depth}\PY{p}{(}\PY{p}{)}\PY{p}{)}
\PY{n+nb}{print}\PY{p}{(}\PY{l+s+s2}{\PYZdq{}}\PY{l+s+s2}{Número de folhas:}\PY{l+s+s2}{\PYZdq{}}\PY{p}{,} \PY{n}{tree\PYZus{}entropy}\PY{o}{.}\PY{n}{get\PYZus{}n\PYZus{}leaves}\PY{p}{(}\PY{p}{)}\PY{p}{)}
\end{Verbatim}
\end{tcolorbox}

    \begin{Verbatim}[commandchars=\\\{\}]
Acurácia: 0.8805662046593925
Precisão: 0.4892367906066536
Recall: 0.4725897920604915
F1: 0.4807692307692308
Profundidade: 36
Número de folhas: 2358
    \end{Verbatim}

    \begin{tcolorbox}[breakable, size=fbox, boxrule=1pt, pad at break*=1mm,colback=cellbackground, colframe=cellborder]
\prompt{In}{incolor}{14}{\boxspacing}
\begin{Verbatim}[commandchars=\\\{\}]
\PY{n}{plt}\PY{o}{.}\PY{n}{figure}\PY{p}{(}\PY{n}{figsize}\PY{o}{=}\PY{p}{(}\PY{l+m+mi}{24}\PY{p}{,} \PY{l+m+mi}{12}\PY{p}{)}\PY{p}{)}
\PY{n}{plot\PYZus{}tree}\PY{p}{(}
    \PY{n}{tree\PYZus{}entropy}\PY{p}{,}
    \PY{n}{max\PYZus{}depth}\PY{o}{=}\PY{l+m+mi}{3}\PY{p}{,}
    \PY{n}{feature\PYZus{}names}\PY{o}{=}\PY{n}{feature\PYZus{}names}\PY{p}{,}
    \PY{n}{class\PYZus{}names}\PY{o}{=}\PY{p}{[}\PY{l+s+s2}{\PYZdq{}}\PY{l+s+s2}{no}\PY{l+s+s2}{\PYZdq{}}\PY{p}{,} \PY{l+s+s2}{\PYZdq{}}\PY{l+s+s2}{yes}\PY{l+s+s2}{\PYZdq{}}\PY{p}{]}\PY{p}{,}
    \PY{n}{filled}\PY{o}{=}\PY{k+kc}{True}\PY{p}{,}
    \PY{n}{fontsize}\PY{o}{=}\PY{l+m+mi}{8}
\PY{p}{)}
\PY{n}{plt}\PY{o}{.}\PY{n}{show}\PY{p}{(}\PY{p}{)}
\end{Verbatim}
\end{tcolorbox}

    \begin{center}
    \adjustimage{max size={0.9\linewidth}{0.9\paperheight}}{output_32_0.png}
    \end{center}
    { \hspace*{\fill} \\}
    
    \subsubsection{Questão 6 --- Gini ×
Entropy}\label{questuxe3o-6-gini-entropy}

Preencha a tabela:

{\def\LTcaptype{none} % do not increment counter
\begin{longtable}[]{@{}lrr@{}}
\toprule\noalign{}
Característica & Gini & Entropy \\
\midrule\noalign{}
\endhead
\bottomrule\noalign{}
\endlastfoot
Acurácia & 0.8731 & 0.8806 \\
Precisão & 0.4593 & 0.4892 \\
Recall & 0.4770 & 0.4726 \\
F1 & 0.4680 & 0.4808 \\
Profundidade & 36 & 36 \\
Número de folhas & 2562 & 2358 \\
Atributo da raiz & duration \textless= 518.5 & duration \textless=
452.5 \\
\end{longtable}
}

Depois responda: 1. As árvores são iguais? 2. O atributo escolhido na
raiz é o mesmo? 3. A ordem dos atributos é a mesma? 4. Qual árvore é
mais complexa? 5. Qual apresentou melhor desempenho no teste? 6. Uma
diferença pequena de desempenho permite afirmar que um critério é sempre
melhor que o outro?

\textbf{Resposta:}

    \begin{enumerate}
\def\labelenumi{\arabic{enumi}.}
\tightlist
\item
  Não, elas têm número de folhas e divisões diferentes.
\item
  Sim, ambas escolheram duration na raiz.
\item
  Não é idêntica nos níveis mais profundos.
\item
  A árvore com critério Gini é mais complexa (2562 folhas contra 2358 de
  Entropy).
\item
  Entropy teve desempenho um pouco melhor no teste (acurácia 88,06\% vs
  87,31\% e F1 0,4808 vs 0,4680).
\item
  Não, a diferença é pequena e na prática os dois critérios são muito
  semelhantes. O controle de profundidade e poda influenciam muito mais
  o resultado final.
\end{enumerate}

    \section{Parte II --- Parada
antecipada}\label{parte-ii-parada-antecipada}

\subsection{10. Investigando o crescimento da
árvore}\label{investigando-o-crescimento-da-uxe1rvore}

Uma das características das Árvores de Decisão é que sua complexidade
pode ser controlada por meio do parâmetro \texttt{max\_depth}.

Até o momento, foram construídos modelos utilizando uma única
configuração de profundidade. Agora, você deverá investigar \textbf{como
o desempenho e a complexidade da árvore variam à medida que sua
profundidade máxima aumenta}.

Considere os seguintes valores para \texttt{max\_depth}:

\begin{Shaded}
\begin{Highlighting}[]
\NormalTok{[}\DecValTok{2}\NormalTok{, }\DecValTok{3}\NormalTok{, }\DecValTok{4}\NormalTok{, }\DecValTok{5}\NormalTok{, }\DecValTok{6}\NormalTok{, }\DecValTok{8}\NormalTok{, }\DecValTok{10}\NormalTok{, }\DecValTok{15}\NormalTok{, }\DecValTok{20}\NormalTok{, }\VariableTok{None}\NormalTok{]}
\end{Highlighting}
\end{Shaded}

O valor \texttt{None} representa uma árvore sem uma restrição explícita
de profundidade.

\paragraph{Tarefa}\label{tarefa}

Escreva um código que, para \textbf{cada valor de \texttt{max\_depth}}:

\begin{enumerate}
\def\labelenumi{\arabic{enumi}.}
\tightlist
\item
  crie uma nova \texttt{DecisionTreeClassifier}, utilizando o critério
  \texttt{gini};
\item
  treine o modelo utilizando \texttt{X\_train} e \texttt{y\_train};
\item
  realize previsões tanto no conjunto de treinamento quanto no conjunto
  de teste;
\item
  calcule a \textbf{acurácia no treinamento};
\item
  calcule a \textbf{acurácia no teste};
\item
  calcule o \textbf{F1-score no teste};
\item
  obtenha a profundidade efetivamente alcançada pela árvore;
\item
  obtenha o número de folhas da árvore;
\item
  armazene todos esses resultados para posterior comparação.
\end{enumerate}

Ao final, organize os resultados em um \texttt{DataFrame} chamado
\texttt{results\_df}.

O \texttt{DataFrame} deverá permitir comparar, para cada profundidade,
pelo menos as seguintes informações:

{\def\LTcaptype{none} % do not increment counter
\begin{longtable}[]{@{}ll@{}}
\toprule\noalign{}
Informação & Descrição \\
\midrule\noalign{}
\endhead
\bottomrule\noalign{}
\endlastfoot
\texttt{max\_depth} & Limite de profundidade utilizado \\
\texttt{train\_accuracy} & Acurácia no treinamento \\
\texttt{test\_accuracy} & Acurácia no teste \\
\texttt{test\_f1} & F1-score no teste \\
\texttt{tree\_depth} & Profundidade efetivamente atingida \\
\texttt{leaves} & Número de folhas \\
\end{longtable}
}

\textbf{Restrição:} não crie manualmente um modelo para cada
profundidade. Utilize uma estrutura de repetição para executar o
experimento automaticamente.

    \begin{tcolorbox}[breakable, size=fbox, boxrule=1pt, pad at break*=1mm,colback=cellbackground, colframe=cellborder]
\prompt{In}{incolor}{15}{\boxspacing}
\begin{Verbatim}[commandchars=\\\{\}]
\PY{n}{depth\PYZus{}list} \PY{o}{=} \PY{p}{[}\PY{l+m+mi}{2}\PY{p}{,} \PY{l+m+mi}{3}\PY{p}{,} \PY{l+m+mi}{4}\PY{p}{,} \PY{l+m+mi}{5}\PY{p}{,} \PY{l+m+mi}{6}\PY{p}{,} \PY{l+m+mi}{8}\PY{p}{,} \PY{l+m+mi}{10}\PY{p}{,} \PY{l+m+mi}{15}\PY{p}{,} \PY{l+m+mi}{20}\PY{p}{,} \PY{k+kc}{None}\PY{p}{]}
\PY{n}{results} \PY{o}{=} \PY{p}{[}\PY{p}{]}

\PY{k}{for} \PY{n}{max\PYZus{}d} \PY{o+ow}{in} \PY{n}{depth\PYZus{}list}\PY{p}{:}
    \PY{n}{clf} \PY{o}{=} \PY{n}{DecisionTreeClassifier}\PY{p}{(}\PY{n}{criterion}\PY{o}{=}\PY{l+s+s1}{\PYZsq{}}\PY{l+s+s1}{gini}\PY{l+s+s1}{\PYZsq{}}\PY{p}{,} \PY{n}{max\PYZus{}depth}\PY{o}{=}\PY{n}{max\PYZus{}d}\PY{p}{,} \PY{n}{random\PYZus{}state}\PY{o}{=}\PY{l+m+mi}{42}\PY{p}{)}
    \PY{n}{clf}\PY{o}{.}\PY{n}{fit}\PY{p}{(}\PY{n}{X\PYZus{}train}\PY{p}{,} \PY{n}{y\PYZus{}train}\PY{p}{)}
    
    \PY{n}{y\PYZus{}pred\PYZus{}tr} \PY{o}{=} \PY{n}{clf}\PY{o}{.}\PY{n}{predict}\PY{p}{(}\PY{n}{X\PYZus{}train}\PY{p}{)}
    \PY{n}{y\PYZus{}pred\PYZus{}te} \PY{o}{=} \PY{n}{clf}\PY{o}{.}\PY{n}{predict}\PY{p}{(}\PY{n}{X\PYZus{}test}\PY{p}{)}
    
    \PY{n}{results}\PY{o}{.}\PY{n}{append}\PY{p}{(}\PY{p}{\PYZob{}}
        \PY{l+s+s1}{\PYZsq{}}\PY{l+s+s1}{max\PYZus{}depth}\PY{l+s+s1}{\PYZsq{}}\PY{p}{:} \PY{n+nb}{str}\PY{p}{(}\PY{n}{max\PYZus{}d}\PY{p}{)}\PY{p}{,}
        \PY{l+s+s1}{\PYZsq{}}\PY{l+s+s1}{train\PYZus{}accuracy}\PY{l+s+s1}{\PYZsq{}}\PY{p}{:} \PY{n}{accuracy\PYZus{}score}\PY{p}{(}\PY{n}{y\PYZus{}train}\PY{p}{,} \PY{n}{y\PYZus{}pred\PYZus{}tr}\PY{p}{)}\PY{p}{,}
        \PY{l+s+s1}{\PYZsq{}}\PY{l+s+s1}{test\PYZus{}accuracy}\PY{l+s+s1}{\PYZsq{}}\PY{p}{:} \PY{n}{accuracy\PYZus{}score}\PY{p}{(}\PY{n}{y\PYZus{}test}\PY{p}{,} \PY{n}{y\PYZus{}pred\PYZus{}te}\PY{p}{)}\PY{p}{,}
        \PY{l+s+s1}{\PYZsq{}}\PY{l+s+s1}{test\PYZus{}f1}\PY{l+s+s1}{\PYZsq{}}\PY{p}{:} \PY{n}{f1\PYZus{}score}\PY{p}{(}\PY{n}{y\PYZus{}test}\PY{p}{,} \PY{n}{y\PYZus{}pred\PYZus{}te}\PY{p}{)}\PY{p}{,}
        \PY{l+s+s1}{\PYZsq{}}\PY{l+s+s1}{tree\PYZus{}depth}\PY{l+s+s1}{\PYZsq{}}\PY{p}{:} \PY{n}{clf}\PY{o}{.}\PY{n}{get\PYZus{}depth}\PY{p}{(}\PY{p}{)}\PY{p}{,}
        \PY{l+s+s1}{\PYZsq{}}\PY{l+s+s1}{leaves}\PY{l+s+s1}{\PYZsq{}}\PY{p}{:} \PY{n}{clf}\PY{o}{.}\PY{n}{get\PYZus{}n\PYZus{}leaves}\PY{p}{(}\PY{p}{)}
    \PY{p}{\PYZcb{}}\PY{p}{)}

\PY{n}{results\PYZus{}df} \PY{o}{=} \PY{n}{pd}\PY{o}{.}\PY{n}{DataFrame}\PY{p}{(}\PY{n}{results}\PY{p}{)}
\PY{n}{results\PYZus{}df}
\end{Verbatim}
\end{tcolorbox}

            \begin{tcolorbox}[breakable, size=fbox, boxrule=.5pt, pad at break*=1mm, opacityfill=0]
\prompt{Out}{outcolor}{15}{\boxspacing}
\begin{Verbatim}[commandchars=\\\{\}]
  max\_depth  train\_accuracy  test\_accuracy   test\_f1  tree\_depth  leaves
0         2        0.896072       0.897007  0.436466           2       4
1         3        0.901476       0.900914  0.446914           3       8
2         4        0.902898       0.899513  0.412753           4      16
3         5        0.905204       0.900767  0.451060           5      31
4         6        0.908933       0.902094  0.473434           6      61
5         8        0.919297       0.902757  0.479684           8     191
6        10        0.931020       0.899808  0.486589          10     432
7        15        0.963914       0.892215  0.497940          15    1274
8        20        0.987708       0.883663  0.485323          20    2036
9      None        1.000000       0.873120  0.468006          36    2562
\end{Verbatim}
\end{tcolorbox}
        
    Usando o dataframe criado, observe o efeito da profundidade sobre o
desempenho.

    \begin{tcolorbox}[breakable, size=fbox, boxrule=1pt, pad at break*=1mm,colback=cellbackground, colframe=cellborder]
\prompt{In}{incolor}{16}{\boxspacing}
\begin{Verbatim}[commandchars=\\\{\}]
\PY{n}{plt}\PY{o}{.}\PY{n}{figure}\PY{p}{(}\PY{n}{figsize}\PY{o}{=}\PY{p}{(}\PY{l+m+mi}{10}\PY{p}{,} \PY{l+m+mi}{6}\PY{p}{)}\PY{p}{)}
\PY{n}{plt}\PY{o}{.}\PY{n}{plot}\PY{p}{(}\PY{n}{results\PYZus{}df}\PY{p}{[}\PY{l+s+s2}{\PYZdq{}}\PY{l+s+s2}{tree\PYZus{}depth}\PY{l+s+s2}{\PYZdq{}}\PY{p}{]}\PY{p}{,} \PY{n}{results\PYZus{}df}\PY{p}{[}\PY{l+s+s2}{\PYZdq{}}\PY{l+s+s2}{train\PYZus{}accuracy}\PY{l+s+s2}{\PYZdq{}}\PY{p}{]}\PY{p}{,}
         \PY{n}{marker}\PY{o}{=}\PY{l+s+s2}{\PYZdq{}}\PY{l+s+s2}{o}\PY{l+s+s2}{\PYZdq{}}\PY{p}{,} \PY{n}{label}\PY{o}{=}\PY{l+s+s2}{\PYZdq{}}\PY{l+s+s2}{Treinamento}\PY{l+s+s2}{\PYZdq{}}\PY{p}{)}
\PY{n}{plt}\PY{o}{.}\PY{n}{plot}\PY{p}{(}\PY{n}{results\PYZus{}df}\PY{p}{[}\PY{l+s+s2}{\PYZdq{}}\PY{l+s+s2}{tree\PYZus{}depth}\PY{l+s+s2}{\PYZdq{}}\PY{p}{]}\PY{p}{,} \PY{n}{results\PYZus{}df}\PY{p}{[}\PY{l+s+s2}{\PYZdq{}}\PY{l+s+s2}{test\PYZus{}accuracy}\PY{l+s+s2}{\PYZdq{}}\PY{p}{]}\PY{p}{,}
         \PY{n}{marker}\PY{o}{=}\PY{l+s+s2}{\PYZdq{}}\PY{l+s+s2}{o}\PY{l+s+s2}{\PYZdq{}}\PY{p}{,} \PY{n}{label}\PY{o}{=}\PY{l+s+s2}{\PYZdq{}}\PY{l+s+s2}{Teste}\PY{l+s+s2}{\PYZdq{}}\PY{p}{)}
\PY{n}{plt}\PY{o}{.}\PY{n}{xlabel}\PY{p}{(}\PY{l+s+s2}{\PYZdq{}}\PY{l+s+s2}{Profundidade da árvore}\PY{l+s+s2}{\PYZdq{}}\PY{p}{)}
\PY{n}{plt}\PY{o}{.}\PY{n}{ylabel}\PY{p}{(}\PY{l+s+s2}{\PYZdq{}}\PY{l+s+s2}{Acurácia}\PY{l+s+s2}{\PYZdq{}}\PY{p}{)}
\PY{n}{plt}\PY{o}{.}\PY{n}{title}\PY{p}{(}\PY{l+s+s2}{\PYZdq{}}\PY{l+s+s2}{Efeito da profundidade sobre o desempenho}\PY{l+s+s2}{\PYZdq{}}\PY{p}{)}
\PY{n}{plt}\PY{o}{.}\PY{n}{legend}\PY{p}{(}\PY{p}{)}
\PY{n}{plt}\PY{o}{.}\PY{n}{grid}\PY{p}{(}\PY{p}{)}
\PY{n}{plt}\PY{o}{.}\PY{n}{show}\PY{p}{(}\PY{p}{)}
\end{Verbatim}
\end{tcolorbox}

    \begin{center}
    \adjustimage{max size={0.9\linewidth}{0.9\paperheight}}{output_38_0.png}
    \end{center}
    { \hspace*{\fill} \\}
    
    \subsubsection{Questão 7 --- Efeito da
profundidade}\label{questuxe3o-7-efeito-da-profundidade}

\begin{enumerate}
\def\labelenumi{\arabic{enumi}.}
\tightlist
\item
  O que acontece com a acurácia de treinamento quando a árvore cresce?
\item
  O comportamento da acurácia de teste é igual?
\item
  Existe uma profundidade a partir da qual aumentar a complexidade deixa
  de produzir benefícios claros?
\item
  Em qual profundidade você escolheria trabalhar?
\item
  Justifique sua escolha considerando \textbf{desempenho e
  complexidade}.
\end{enumerate}

\textbf{Resposta:}

    Resposta:

\begin{enumerate}
\def\labelenumi{\arabic{enumi}.}
\tightlist
\item
  A acurácia de treinamento cresce continuamente até atingir 100\%
  (1.000) no modelo sem restrição.
\item
  Não. A acurácia de teste sobe até profundidades entre 5 e 8 e depois
  começa a cair (de \textasciitilde90,3\% para 87,3\%), indicando
  overfitting.
\item
  Sim, a partir de max\_depth = 5 ou 6 o ganho no teste se estabiliza e
  depois piora.
\item
  Escolheria max\_depth = 5 (ou 6).
\item
  Com profundidade 5 o modelo atinge 90,08\% de acurácia no teste com
  apenas 31 folhas. É um modelo simples, fácil de interpretar e que
  evita o sobreajuste da árvore sem limite.
\end{enumerate}

    \subsection{11. Construindo o modelo com parada
antecipada}\label{construindo-o-modelo-com-parada-antecipada}

Escolha uma profundidade com base na análise anterior.

\begin{quote}
Altere SUA\_PROFUNDIDADE para o valor escolhido por você.
\end{quote}

    \begin{tcolorbox}[breakable, size=fbox, boxrule=1pt, pad at break*=1mm,colback=cellbackground, colframe=cellborder]
\prompt{In}{incolor}{17}{\boxspacing}
\begin{Verbatim}[commandchars=\\\{\}]
\PY{n}{SUA\PYZus{}PROFUNDIDADE} \PY{o}{=} \PY{l+m+mi}{6}  \PY{c+c1}{\PYZsh{} altere após sua análise}

\PY{n}{tree\PYZus{}pruned} \PY{o}{=} \PY{n}{DecisionTreeClassifier}\PY{p}{(}
    \PY{n}{criterion}\PY{o}{=}\PY{l+s+s2}{\PYZdq{}}\PY{l+s+s2}{gini}\PY{l+s+s2}{\PYZdq{}}\PY{p}{,}
    \PY{n}{max\PYZus{}depth}\PY{o}{=}\PY{n}{SUA\PYZus{}PROFUNDIDADE}\PY{p}{,}
    \PY{n}{random\PYZus{}state}\PY{o}{=}\PY{l+m+mi}{42}
\PY{p}{)}

\PY{n}{tree\PYZus{}pruned}\PY{o}{.}\PY{n}{fit}\PY{p}{(}\PY{n}{X\PYZus{}train}\PY{p}{,} \PY{n}{y\PYZus{}train}\PY{p}{)}
\PY{n}{y\PYZus{}pred\PYZus{}pruned} \PY{o}{=} \PY{n}{tree\PYZus{}pruned}\PY{o}{.}\PY{n}{predict}\PY{p}{(}\PY{n}{X\PYZus{}test}\PY{p}{)}

\PY{n+nb}{print}\PY{p}{(}\PY{l+s+s2}{\PYZdq{}}\PY{l+s+s2}{Acurácia:}\PY{l+s+s2}{\PYZdq{}}\PY{p}{,} \PY{n}{accuracy\PYZus{}score}\PY{p}{(}\PY{n}{y\PYZus{}test}\PY{p}{,} \PY{n}{y\PYZus{}pred\PYZus{}pruned}\PY{p}{)}\PY{p}{)}
\PY{n+nb}{print}\PY{p}{(}\PY{l+s+s2}{\PYZdq{}}\PY{l+s+s2}{Precisão:}\PY{l+s+s2}{\PYZdq{}}\PY{p}{,} \PY{n}{precision\PYZus{}score}\PY{p}{(}\PY{n}{y\PYZus{}test}\PY{p}{,} \PY{n}{y\PYZus{}pred\PYZus{}pruned}\PY{p}{)}\PY{p}{)}
\PY{n+nb}{print}\PY{p}{(}\PY{l+s+s2}{\PYZdq{}}\PY{l+s+s2}{Recall:}\PY{l+s+s2}{\PYZdq{}}\PY{p}{,} \PY{n}{recall\PYZus{}score}\PY{p}{(}\PY{n}{y\PYZus{}test}\PY{p}{,} \PY{n}{y\PYZus{}pred\PYZus{}pruned}\PY{p}{)}\PY{p}{)}
\PY{n+nb}{print}\PY{p}{(}\PY{l+s+s2}{\PYZdq{}}\PY{l+s+s2}{F1:}\PY{l+s+s2}{\PYZdq{}}\PY{p}{,} \PY{n}{f1\PYZus{}score}\PY{p}{(}\PY{n}{y\PYZus{}test}\PY{p}{,} \PY{n}{y\PYZus{}pred\PYZus{}pruned}\PY{p}{)}\PY{p}{)}
\PY{n+nb}{print}\PY{p}{(}\PY{l+s+s2}{\PYZdq{}}\PY{l+s+s2}{Profundidade:}\PY{l+s+s2}{\PYZdq{}}\PY{p}{,} \PY{n}{tree\PYZus{}pruned}\PY{o}{.}\PY{n}{get\PYZus{}depth}\PY{p}{(}\PY{p}{)}\PY{p}{)}
\PY{n+nb}{print}\PY{p}{(}\PY{l+s+s2}{\PYZdq{}}\PY{l+s+s2}{Número de folhas:}\PY{l+s+s2}{\PYZdq{}}\PY{p}{,} \PY{n}{tree\PYZus{}pruned}\PY{o}{.}\PY{n}{get\PYZus{}n\PYZus{}leaves}\PY{p}{(}\PY{p}{)}\PY{p}{)}
\end{Verbatim}
\end{tcolorbox}

    \begin{Verbatim}[commandchars=\\\{\}]
Acurácia: 0.9020937776467118
Precisão: 0.6385026737967915
Recall: 0.3761814744801512
F1: 0.4734337827121332
Profundidade: 6
Número de folhas: 61
    \end{Verbatim}

    \begin{tcolorbox}[breakable, size=fbox, boxrule=1pt, pad at break*=1mm,colback=cellbackground, colframe=cellborder]
\prompt{In}{incolor}{18}{\boxspacing}
\begin{Verbatim}[commandchars=\\\{\}]
\PY{n}{plt}\PY{o}{.}\PY{n}{figure}\PY{p}{(}\PY{n}{figsize}\PY{o}{=}\PY{p}{(}\PY{l+m+mi}{24}\PY{p}{,} \PY{l+m+mi}{12}\PY{p}{)}\PY{p}{)}
\PY{n}{plot\PYZus{}tree}\PY{p}{(}
    \PY{n}{tree\PYZus{}pruned}\PY{p}{,}
    \PY{n}{feature\PYZus{}names}\PY{o}{=}\PY{n}{feature\PYZus{}names}\PY{p}{,}
    \PY{n}{class\PYZus{}names}\PY{o}{=}\PY{p}{[}\PY{l+s+s2}{\PYZdq{}}\PY{l+s+s2}{no}\PY{l+s+s2}{\PYZdq{}}\PY{p}{,} \PY{l+s+s2}{\PYZdq{}}\PY{l+s+s2}{yes}\PY{l+s+s2}{\PYZdq{}}\PY{p}{]}\PY{p}{,}
    \PY{n}{filled}\PY{o}{=}\PY{k+kc}{True}\PY{p}{,}
    \PY{n}{fontsize}\PY{o}{=}\PY{l+m+mi}{8}
\PY{p}{)}
\PY{n}{plt}\PY{o}{.}\PY{n}{show}\PY{p}{(}\PY{p}{)}
\end{Verbatim}
\end{tcolorbox}

    \begin{center}
    \adjustimage{max size={0.9\linewidth}{0.9\paperheight}}{output_43_0.png}
    \end{center}
    { \hspace*{\fill} \\}
    
    \section{Parte III --- Análise dos
atributos}\label{parte-iii-anuxe1lise-dos-atributos}

\subsection{12. Importância dos
atributos}\label{importuxe2ncia-dos-atributos}

O código a seguir calcula as importâncias dos atributos e as ordena.

    \begin{tcolorbox}[breakable, size=fbox, boxrule=1pt, pad at break*=1mm,colback=cellbackground, colframe=cellborder]
\prompt{In}{incolor}{19}{\boxspacing}
\begin{Verbatim}[commandchars=\\\{\}]
\PY{n}{importance} \PY{o}{=} \PY{n}{pd}\PY{o}{.}\PY{n}{DataFrame}\PY{p}{(}\PY{p}{\PYZob{}}
    \PY{l+s+s2}{\PYZdq{}}\PY{l+s+s2}{feature}\PY{l+s+s2}{\PYZdq{}}\PY{p}{:} \PY{n}{feature\PYZus{}names}\PY{p}{,}
    \PY{l+s+s2}{\PYZdq{}}\PY{l+s+s2}{importance}\PY{l+s+s2}{\PYZdq{}}\PY{p}{:} \PY{n}{tree\PYZus{}gini}\PY{o}{.}\PY{n}{feature\PYZus{}importances\PYZus{}}
\PY{p}{\PYZcb{}}\PY{p}{)}\PY{o}{.}\PY{n}{sort\PYZus{}values}\PY{p}{(}\PY{l+s+s2}{\PYZdq{}}\PY{l+s+s2}{importance}\PY{l+s+s2}{\PYZdq{}}\PY{p}{,} \PY{n}{ascending}\PY{o}{=}\PY{k+kc}{False}\PY{p}{)}

\PY{n}{importance}\PY{o}{.}\PY{n}{head}\PY{p}{(}\PY{l+m+mi}{15}\PY{p}{)}
\end{Verbatim}
\end{tcolorbox}

            \begin{tcolorbox}[breakable, size=fbox, boxrule=.5pt, pad at break*=1mm, opacityfill=0]
\prompt{Out}{outcolor}{19}{\boxspacing}
\begin{Verbatim}[commandchars=\\\{\}]
                  feature  importance
47    remainder\_\_duration    0.279931
45     remainder\_\_balance    0.095818
44         remainder\_\_age    0.093600
46         remainder\_\_day    0.089978
42  cat\_\_poutcome\_success    0.089024
49       remainder\_\_pdays    0.039845
48    remainder\_\_campaign    0.031905
21        cat\_\_housing\_no    0.017494
28         cat\_\_month\_apr    0.015839
35         cat\_\_month\_mar    0.014490
34         cat\_\_month\_jun    0.014359
50    remainder\_\_previous    0.012954
38         cat\_\_month\_oct    0.011874
27   cat\_\_contact\_unknown    0.010367
36         cat\_\_month\_may    0.009264
\end{Verbatim}
\end{tcolorbox}
        
    \begin{tcolorbox}[breakable, size=fbox, boxrule=1pt, pad at break*=1mm,colback=cellbackground, colframe=cellborder]
\prompt{In}{incolor}{20}{\boxspacing}
\begin{Verbatim}[commandchars=\\\{\}]
\PY{n}{importance}\PY{o}{.}\PY{n}{head}\PY{p}{(}\PY{l+m+mi}{15}\PY{p}{)}\PY{o}{.}\PY{n}{plot}\PY{p}{(}
    \PY{n}{x}\PY{o}{=}\PY{l+s+s2}{\PYZdq{}}\PY{l+s+s2}{feature}\PY{l+s+s2}{\PYZdq{}}\PY{p}{,}
    \PY{n}{y}\PY{o}{=}\PY{l+s+s2}{\PYZdq{}}\PY{l+s+s2}{importance}\PY{l+s+s2}{\PYZdq{}}\PY{p}{,}
    \PY{n}{kind}\PY{o}{=}\PY{l+s+s2}{\PYZdq{}}\PY{l+s+s2}{barh}\PY{l+s+s2}{\PYZdq{}}\PY{p}{,}
    \PY{n}{figsize}\PY{o}{=}\PY{p}{(}\PY{l+m+mi}{10}\PY{p}{,} \PY{l+m+mi}{6}\PY{p}{)}
\PY{p}{)}
\PY{n}{plt}\PY{o}{.}\PY{n}{xlabel}\PY{p}{(}\PY{l+s+s2}{\PYZdq{}}\PY{l+s+s2}{Importância}\PY{l+s+s2}{\PYZdq{}}\PY{p}{)}
\PY{n}{plt}\PY{o}{.}\PY{n}{ylabel}\PY{p}{(}\PY{l+s+s2}{\PYZdq{}}\PY{l+s+s2}{Atributo}\PY{l+s+s2}{\PYZdq{}}\PY{p}{)}
\PY{n}{plt}\PY{o}{.}\PY{n}{title}\PY{p}{(}\PY{l+s+s2}{\PYZdq{}}\PY{l+s+s2}{Principais atributos segundo a árvore}\PY{l+s+s2}{\PYZdq{}}\PY{p}{)}
\PY{n}{plt}\PY{o}{.}\PY{n}{show}\PY{p}{(}\PY{p}{)}
\end{Verbatim}
\end{tcolorbox}

    \begin{center}
    \adjustimage{max size={0.9\linewidth}{0.9\paperheight}}{output_46_0.png}
    \end{center}
    { \hspace*{\fill} \\}
    
    \subsubsection{Questão 8 --- Atributos
escolhidos}\label{questuxe3o-8-atributos-escolhidos}

\begin{enumerate}
\def\labelenumi{\arabic{enumi}.}
\tightlist
\item
  Quais são os cinco atributos com maior importância?
\item
  O atributo utilizado na raiz está entre os mais importantes?
\item
  Um atributo que aparece em um nível profundo da árvore necessariamente
  é menos importante?
\item
  Explique a diferença entre:

  \begin{itemize}
  \tightlist
  \item
    atributo escolhido para o primeiro particionamento;
  \item
    atributo com maior importância no modelo.
  \end{itemize}
\end{enumerate}

\textbf{Resposta:}

    \begin{enumerate}
\def\labelenumi{\arabic{enumi}.}
\tightlist
\item
  Os 5 atributos mais importantes são:

  \begin{itemize}
  \tightlist
  \item
    duration (\textasciitilde28,0\%)
  \item
    balance (\textasciitilde9,6\%)
  \item
    age (\textasciitilde9,4\%)
  \item
    day (\textasciitilde9,0\%)
  \item
    poutcome\_success (\textasciitilde8,9\%)
  \end{itemize}
\item
  Sim, duration é o atributo mais importante de todo o modelo.
\item
  Não necessariamente, mas no nível profundo cada divisão afeta menos
  amostras. Contudo, se uma variável for usada muitas vezes em vários
  nós, ela pode somar bastante importância.
\item
  \begin{itemize}
  \tightlist
  \item
    \textbf{Atributo da raiz:} é o melhor atributo na primeira divisão
    inicial (decisão local no topo da árvore).
  \item
    \textbf{Atributo com maior importância:} é a soma de todas as
    reduções de impureza que a variável gerou em todas as divisões da
    árvore (métrica global).
  \end{itemize}
\end{enumerate}

    \section{Parte IV --- Comparação
final}\label{parte-iv-comparauxe7uxe3o-final}

\subsection{13. Comparando os modelos}\label{comparando-os-modelos}

Compare: 1. Gini sem restrição; 2. Entropy sem restrição; 3. Gini com
parada antecipada.

    \begin{tcolorbox}[breakable, size=fbox, boxrule=1pt, pad at break*=1mm,colback=cellbackground, colframe=cellborder]
\prompt{In}{incolor}{21}{\boxspacing}
\begin{Verbatim}[commandchars=\\\{\}]
\PY{n}{comparison} \PY{o}{=} \PY{n}{pd}\PY{o}{.}\PY{n}{DataFrame}\PY{p}{(}\PY{p}{\PYZob{}}
    \PY{l+s+s2}{\PYZdq{}}\PY{l+s+s2}{Modelo}\PY{l+s+s2}{\PYZdq{}}\PY{p}{:} \PY{p}{[}
        \PY{l+s+s2}{\PYZdq{}}\PY{l+s+s2}{Gini \PYZhy{} sem restrição}\PY{l+s+s2}{\PYZdq{}}\PY{p}{,}
        \PY{l+s+s2}{\PYZdq{}}\PY{l+s+s2}{Entropy \PYZhy{} sem restrição}\PY{l+s+s2}{\PYZdq{}}\PY{p}{,}
        \PY{l+s+s2}{\PYZdq{}}\PY{l+s+s2}{Gini \PYZhy{} parada antecipada}\PY{l+s+s2}{\PYZdq{}}
    \PY{p}{]}\PY{p}{,}
    \PY{l+s+s2}{\PYZdq{}}\PY{l+s+s2}{Acurácia}\PY{l+s+s2}{\PYZdq{}}\PY{p}{:} \PY{p}{[}
        \PY{n}{accuracy\PYZus{}score}\PY{p}{(}\PY{n}{y\PYZus{}test}\PY{p}{,} \PY{n}{tree\PYZus{}gini}\PY{o}{.}\PY{n}{predict}\PY{p}{(}\PY{n}{X\PYZus{}test}\PY{p}{)}\PY{p}{)}\PY{p}{,}
        \PY{n}{accuracy\PYZus{}score}\PY{p}{(}\PY{n}{y\PYZus{}test}\PY{p}{,} \PY{n}{tree\PYZus{}entropy}\PY{o}{.}\PY{n}{predict}\PY{p}{(}\PY{n}{X\PYZus{}test}\PY{p}{)}\PY{p}{)}\PY{p}{,}
        \PY{n}{accuracy\PYZus{}score}\PY{p}{(}\PY{n}{y\PYZus{}test}\PY{p}{,} \PY{n}{tree\PYZus{}pruned}\PY{o}{.}\PY{n}{predict}\PY{p}{(}\PY{n}{X\PYZus{}test}\PY{p}{)}\PY{p}{)}
    \PY{p}{]}\PY{p}{,}
    \PY{l+s+s2}{\PYZdq{}}\PY{l+s+s2}{Precisão}\PY{l+s+s2}{\PYZdq{}}\PY{p}{:} \PY{p}{[}
        \PY{n}{precision\PYZus{}score}\PY{p}{(}\PY{n}{y\PYZus{}test}\PY{p}{,} \PY{n}{tree\PYZus{}gini}\PY{o}{.}\PY{n}{predict}\PY{p}{(}\PY{n}{X\PYZus{}test}\PY{p}{)}\PY{p}{)}\PY{p}{,}
        \PY{n}{precision\PYZus{}score}\PY{p}{(}\PY{n}{y\PYZus{}test}\PY{p}{,} \PY{n}{tree\PYZus{}entropy}\PY{o}{.}\PY{n}{predict}\PY{p}{(}\PY{n}{X\PYZus{}test}\PY{p}{)}\PY{p}{)}\PY{p}{,}
        \PY{n}{precision\PYZus{}score}\PY{p}{(}\PY{n}{y\PYZus{}test}\PY{p}{,} \PY{n}{tree\PYZus{}pruned}\PY{o}{.}\PY{n}{predict}\PY{p}{(}\PY{n}{X\PYZus{}test}\PY{p}{)}\PY{p}{)}
    \PY{p}{]}\PY{p}{,}
    \PY{l+s+s2}{\PYZdq{}}\PY{l+s+s2}{Recall}\PY{l+s+s2}{\PYZdq{}}\PY{p}{:} \PY{p}{[}
        \PY{n}{recall\PYZus{}score}\PY{p}{(}\PY{n}{y\PYZus{}test}\PY{p}{,} \PY{n}{tree\PYZus{}gini}\PY{o}{.}\PY{n}{predict}\PY{p}{(}\PY{n}{X\PYZus{}test}\PY{p}{)}\PY{p}{)}\PY{p}{,}
        \PY{n}{recall\PYZus{}score}\PY{p}{(}\PY{n}{y\PYZus{}test}\PY{p}{,} \PY{n}{tree\PYZus{}entropy}\PY{o}{.}\PY{n}{predict}\PY{p}{(}\PY{n}{X\PYZus{}test}\PY{p}{)}\PY{p}{)}\PY{p}{,}
        \PY{n}{recall\PYZus{}score}\PY{p}{(}\PY{n}{y\PYZus{}test}\PY{p}{,} \PY{n}{tree\PYZus{}pruned}\PY{o}{.}\PY{n}{predict}\PY{p}{(}\PY{n}{X\PYZus{}test}\PY{p}{)}\PY{p}{)}
    \PY{p}{]}\PY{p}{,}
    \PY{l+s+s2}{\PYZdq{}}\PY{l+s+s2}{F1}\PY{l+s+s2}{\PYZdq{}}\PY{p}{:} \PY{p}{[}
        \PY{n}{f1\PYZus{}score}\PY{p}{(}\PY{n}{y\PYZus{}test}\PY{p}{,} \PY{n}{tree\PYZus{}gini}\PY{o}{.}\PY{n}{predict}\PY{p}{(}\PY{n}{X\PYZus{}test}\PY{p}{)}\PY{p}{)}\PY{p}{,}
        \PY{n}{f1\PYZus{}score}\PY{p}{(}\PY{n}{y\PYZus{}test}\PY{p}{,} \PY{n}{tree\PYZus{}entropy}\PY{o}{.}\PY{n}{predict}\PY{p}{(}\PY{n}{X\PYZus{}test}\PY{p}{)}\PY{p}{)}\PY{p}{,}
        \PY{n}{f1\PYZus{}score}\PY{p}{(}\PY{n}{y\PYZus{}test}\PY{p}{,} \PY{n}{tree\PYZus{}pruned}\PY{o}{.}\PY{n}{predict}\PY{p}{(}\PY{n}{X\PYZus{}test}\PY{p}{)}\PY{p}{)}
    \PY{p}{]}\PY{p}{,}
    \PY{l+s+s2}{\PYZdq{}}\PY{l+s+s2}{Profundidade}\PY{l+s+s2}{\PYZdq{}}\PY{p}{:} \PY{p}{[}
        \PY{n}{tree\PYZus{}gini}\PY{o}{.}\PY{n}{get\PYZus{}depth}\PY{p}{(}\PY{p}{)}\PY{p}{,}
        \PY{n}{tree\PYZus{}entropy}\PY{o}{.}\PY{n}{get\PYZus{}depth}\PY{p}{(}\PY{p}{)}\PY{p}{,}
        \PY{n}{tree\PYZus{}pruned}\PY{o}{.}\PY{n}{get\PYZus{}depth}\PY{p}{(}\PY{p}{)}
    \PY{p}{]}\PY{p}{,}
    \PY{l+s+s2}{\PYZdq{}}\PY{l+s+s2}{Folhas}\PY{l+s+s2}{\PYZdq{}}\PY{p}{:} \PY{p}{[}
        \PY{n}{tree\PYZus{}gini}\PY{o}{.}\PY{n}{get\PYZus{}n\PYZus{}leaves}\PY{p}{(}\PY{p}{)}\PY{p}{,}
        \PY{n}{tree\PYZus{}entropy}\PY{o}{.}\PY{n}{get\PYZus{}n\PYZus{}leaves}\PY{p}{(}\PY{p}{)}\PY{p}{,}
        \PY{n}{tree\PYZus{}pruned}\PY{o}{.}\PY{n}{get\PYZus{}n\PYZus{}leaves}\PY{p}{(}\PY{p}{)}
    \PY{p}{]}
\PY{p}{\PYZcb{}}\PY{p}{)}

\PY{n}{comparison}
\end{Verbatim}
\end{tcolorbox}

            \begin{tcolorbox}[breakable, size=fbox, boxrule=.5pt, pad at break*=1mm, opacityfill=0]
\prompt{Out}{outcolor}{21}{\boxspacing}
\begin{Verbatim}[commandchars=\\\{\}]
                     Modelo  Acurácia  Precisão    Recall        F1  \textbackslash{}
0      Gini - sem restrição  0.873120  0.459345  0.477001  0.468006
1   Entropy - sem restrição  0.880566  0.489237  0.472590  0.480769
2  Gini - parada antecipada  0.902094  0.638503  0.376181  0.473434

   Profundidade  Folhas
0            36    2562
1            36    2358
2             6      61
\end{Verbatim}
\end{tcolorbox}
        
    \subsubsection{Questão 9 --- Escolha do modelo
final}\label{questuxe3o-9-escolha-do-modelo-final}

Com base nos resultados:

\begin{enumerate}
\def\labelenumi{\arabic{enumi}.}
\tightlist
\item
  Qual modelo apresenta o melhor desempenho?
\item
  Qual modelo apresenta menor complexidade?
\item
  Existe um modelo que ofereça um bom compromisso entre desempenho e
  interpretabilidade?
\item
  Qual modelo você escolheria para uma aplicação real? Justifique.
\end{enumerate}

\textbf{:}

    \begin{enumerate}
\def\labelenumi{\arabic{enumi}.}
\tightlist
\item
  O modelo Gini com parada antecipada (max\_depth=5) teve o melhor
  desempenho no teste (acurácia de 90,08\% e precisão de 63,93\%).
\item
  O modelo com parada antecipada (profundidade 5 e 31 folhas, contra
  profundidade 36 e mais de 2300 folhas dos outros).
\item
  Sim, o modelo com max\_depth=5 oferece o melhor equilíbrio.
\item
  Escolheria a árvore com max\_depth=5 porque tem melhor acurácia no
  teste, evita overfitting e é uma árvore compacta e fácil de
  interpretar como regras de negócio.
\end{enumerate}

    \section{15. Outro mecanismo de parada}\label{outro-mecanismo-de-parada}

Escolha \textbf{um} dos parâmetros abaixo e faça um pequeno experimento:
- \texttt{min\_samples\_split}; - \texttt{min\_samples\_leaf}; -
\texttt{min\_impurity\_decrease}.

Compare-o com o modelo obtido utilizando \texttt{max\_depth}.

Responda: 1. Qual parâmetro você escolheu? 2. Qual é seu efeito? 3. Como
ele alterou a árvore? 4. Como alterou o desempenho? 5. Qual mecanismo de
parada você considera mais adequado para este conjunto de dados?

\textbf{Resposta:}

    \begin{tcolorbox}[breakable, size=fbox, boxrule=1pt, pad at break*=1mm,colback=cellbackground, colframe=cellborder]
\prompt{In}{incolor}{22}{\boxspacing}
\begin{Verbatim}[commandchars=\\\{\}]
\PY{n}{leaf\PYZus{}values} \PY{o}{=} \PY{p}{[}\PY{l+m+mi}{1}\PY{p}{,} \PY{l+m+mi}{5}\PY{p}{,} \PY{l+m+mi}{10}\PY{p}{,} \PY{l+m+mi}{20}\PY{p}{,} \PY{l+m+mi}{50}\PY{p}{,} \PY{l+m+mi}{100}\PY{p}{,} \PY{l+m+mi}{200}\PY{p}{,} \PY{l+m+mi}{500}\PY{p}{]}
\PY{n}{leaf\PYZus{}results} \PY{o}{=} \PY{p}{[}\PY{p}{]}

\PY{k}{for} \PY{n}{min\PYZus{}leaf} \PY{o+ow}{in} \PY{n}{leaf\PYZus{}values}\PY{p}{:}
    \PY{n}{clf\PYZus{}leaf} \PY{o}{=} \PY{n}{DecisionTreeClassifier}\PY{p}{(}\PY{n}{criterion}\PY{o}{=}\PY{l+s+s1}{\PYZsq{}}\PY{l+s+s1}{gini}\PY{l+s+s1}{\PYZsq{}}\PY{p}{,} \PY{n}{min\PYZus{}samples\PYZus{}leaf}\PY{o}{=}\PY{n}{min\PYZus{}leaf}\PY{p}{,} \PY{n}{random\PYZus{}state}\PY{o}{=}\PY{l+m+mi}{42}\PY{p}{)}
    \PY{n}{clf\PYZus{}leaf}\PY{o}{.}\PY{n}{fit}\PY{p}{(}\PY{n}{X\PYZus{}train}\PY{p}{,} \PY{n}{y\PYZus{}train}\PY{p}{)}
    
    \PY{n}{y\PYZus{}pred\PYZus{}te} \PY{o}{=} \PY{n}{clf\PYZus{}leaf}\PY{o}{.}\PY{n}{predict}\PY{p}{(}\PY{n}{X\PYZus{}test}\PY{p}{)}
    \PY{n}{leaf\PYZus{}results}\PY{o}{.}\PY{n}{append}\PY{p}{(}\PY{p}{\PYZob{}}
        \PY{l+s+s1}{\PYZsq{}}\PY{l+s+s1}{min\PYZus{}samples\PYZus{}leaf}\PY{l+s+s1}{\PYZsq{}}\PY{p}{:} \PY{n}{min\PYZus{}leaf}\PY{p}{,}
        \PY{l+s+s1}{\PYZsq{}}\PY{l+s+s1}{test\PYZus{}accuracy}\PY{l+s+s1}{\PYZsq{}}\PY{p}{:} \PY{n}{accuracy\PYZus{}score}\PY{p}{(}\PY{n}{y\PYZus{}test}\PY{p}{,} \PY{n}{y\PYZus{}pred\PYZus{}te}\PY{p}{)}\PY{p}{,}
        \PY{l+s+s1}{\PYZsq{}}\PY{l+s+s1}{test\PYZus{}precision}\PY{l+s+s1}{\PYZsq{}}\PY{p}{:} \PY{n}{precision\PYZus{}score}\PY{p}{(}\PY{n}{y\PYZus{}test}\PY{p}{,} \PY{n}{y\PYZus{}pred\PYZus{}te}\PY{p}{)}\PY{p}{,}
        \PY{l+s+s1}{\PYZsq{}}\PY{l+s+s1}{test\PYZus{}recall}\PY{l+s+s1}{\PYZsq{}}\PY{p}{:} \PY{n}{recall\PYZus{}score}\PY{p}{(}\PY{n}{y\PYZus{}test}\PY{p}{,} \PY{n}{y\PYZus{}pred\PYZus{}te}\PY{p}{)}\PY{p}{,}
        \PY{l+s+s1}{\PYZsq{}}\PY{l+s+s1}{test\PYZus{}f1}\PY{l+s+s1}{\PYZsq{}}\PY{p}{:} \PY{n}{f1\PYZus{}score}\PY{p}{(}\PY{n}{y\PYZus{}test}\PY{p}{,} \PY{n}{y\PYZus{}pred\PYZus{}te}\PY{p}{)}\PY{p}{,}
        \PY{l+s+s1}{\PYZsq{}}\PY{l+s+s1}{tree\PYZus{}depth}\PY{l+s+s1}{\PYZsq{}}\PY{p}{:} \PY{n}{clf\PYZus{}leaf}\PY{o}{.}\PY{n}{get\PYZus{}depth}\PY{p}{(}\PY{p}{)}\PY{p}{,}
        \PY{l+s+s1}{\PYZsq{}}\PY{l+s+s1}{leaves}\PY{l+s+s1}{\PYZsq{}}\PY{p}{:} \PY{n}{clf\PYZus{}leaf}\PY{o}{.}\PY{n}{get\PYZus{}n\PYZus{}leaves}\PY{p}{(}\PY{p}{)}
    \PY{p}{\PYZcb{}}\PY{p}{)}

\PY{n}{leaf\PYZus{}df} \PY{o}{=} \PY{n}{pd}\PY{o}{.}\PY{n}{DataFrame}\PY{p}{(}\PY{n}{leaf\PYZus{}results}\PY{p}{)}
\PY{n}{display}\PY{p}{(}\PY{n}{leaf\PYZus{}df}\PY{p}{)}

\PY{n}{best\PYZus{}leaf\PYZus{}model} \PY{o}{=} \PY{n}{DecisionTreeClassifier}\PY{p}{(}\PY{n}{criterion}\PY{o}{=}\PY{l+s+s1}{\PYZsq{}}\PY{l+s+s1}{gini}\PY{l+s+s1}{\PYZsq{}}\PY{p}{,} \PY{n}{min\PYZus{}samples\PYZus{}leaf}\PY{o}{=}\PY{l+m+mi}{20}\PY{p}{,} \PY{n}{random\PYZus{}state}\PY{o}{=}\PY{l+m+mi}{42}\PY{p}{)}
\PY{n}{best\PYZus{}leaf\PYZus{}model}\PY{o}{.}\PY{n}{fit}\PY{p}{(}\PY{n}{X\PYZus{}train}\PY{p}{,} \PY{n}{y\PYZus{}train}\PY{p}{)}

\PY{n}{final\PYZus{}comp} \PY{o}{=} \PY{n}{pd}\PY{o}{.}\PY{n}{DataFrame}\PY{p}{(}\PY{p}{\PYZob{}}
    \PY{l+s+s1}{\PYZsq{}}\PY{l+s+s1}{Modelo}\PY{l+s+s1}{\PYZsq{}}\PY{p}{:} \PY{p}{[}\PY{l+s+s1}{\PYZsq{}}\PY{l+s+s1}{Sem restrição}\PY{l+s+s1}{\PYZsq{}}\PY{p}{,} \PY{l+s+s1}{\PYZsq{}}\PY{l+s+s1}{max\PYZus{}depth=5}\PY{l+s+s1}{\PYZsq{}}\PY{p}{,} \PY{l+s+s1}{\PYZsq{}}\PY{l+s+s1}{min\PYZus{}samples\PYZus{}leaf=20}\PY{l+s+s1}{\PYZsq{}}\PY{p}{]}\PY{p}{,}
    \PY{l+s+s1}{\PYZsq{}}\PY{l+s+s1}{Acurácia}\PY{l+s+s1}{\PYZsq{}}\PY{p}{:} \PY{p}{[}\PY{n}{accuracy\PYZus{}score}\PY{p}{(}\PY{n}{y\PYZus{}test}\PY{p}{,} \PY{n}{tree\PYZus{}gini}\PY{o}{.}\PY{n}{predict}\PY{p}{(}\PY{n}{X\PYZus{}test}\PY{p}{)}\PY{p}{)}\PY{p}{,} \PY{n}{accuracy\PYZus{}score}\PY{p}{(}\PY{n}{y\PYZus{}test}\PY{p}{,} \PY{n}{tree\PYZus{}pruned}\PY{o}{.}\PY{n}{predict}\PY{p}{(}\PY{n}{X\PYZus{}test}\PY{p}{)}\PY{p}{)}\PY{p}{,} \PY{n}{accuracy\PYZus{}score}\PY{p}{(}\PY{n}{y\PYZus{}test}\PY{p}{,} \PY{n}{best\PYZus{}leaf\PYZus{}model}\PY{o}{.}\PY{n}{predict}\PY{p}{(}\PY{n}{X\PYZus{}test}\PY{p}{)}\PY{p}{)}\PY{p}{]}\PY{p}{,}
    \PY{l+s+s1}{\PYZsq{}}\PY{l+s+s1}{Precisão}\PY{l+s+s1}{\PYZsq{}}\PY{p}{:} \PY{p}{[}\PY{n}{precision\PYZus{}score}\PY{p}{(}\PY{n}{y\PYZus{}test}\PY{p}{,} \PY{n}{tree\PYZus{}gini}\PY{o}{.}\PY{n}{predict}\PY{p}{(}\PY{n}{X\PYZus{}test}\PY{p}{)}\PY{p}{)}\PY{p}{,} \PY{n}{precision\PYZus{}score}\PY{p}{(}\PY{n}{y\PYZus{}test}\PY{p}{,} \PY{n}{tree\PYZus{}pruned}\PY{o}{.}\PY{n}{predict}\PY{p}{(}\PY{n}{X\PYZus{}test}\PY{p}{)}\PY{p}{)}\PY{p}{,} \PY{n}{precision\PYZus{}score}\PY{p}{(}\PY{n}{y\PYZus{}test}\PY{p}{,} \PY{n}{best\PYZus{}leaf\PYZus{}model}\PY{o}{.}\PY{n}{predict}\PY{p}{(}\PY{n}{X\PYZus{}test}\PY{p}{)}\PY{p}{)}\PY{p}{]}\PY{p}{,}
    \PY{l+s+s1}{\PYZsq{}}\PY{l+s+s1}{Recall}\PY{l+s+s1}{\PYZsq{}}\PY{p}{:} \PY{p}{[}\PY{n}{recall\PYZus{}score}\PY{p}{(}\PY{n}{y\PYZus{}test}\PY{p}{,} \PY{n}{tree\PYZus{}gini}\PY{o}{.}\PY{n}{predict}\PY{p}{(}\PY{n}{X\PYZus{}test}\PY{p}{)}\PY{p}{)}\PY{p}{,} \PY{n}{recall\PYZus{}score}\PY{p}{(}\PY{n}{y\PYZus{}test}\PY{p}{,} \PY{n}{tree\PYZus{}pruned}\PY{o}{.}\PY{n}{predict}\PY{p}{(}\PY{n}{X\PYZus{}test}\PY{p}{)}\PY{p}{)}\PY{p}{,} \PY{n}{recall\PYZus{}score}\PY{p}{(}\PY{n}{y\PYZus{}test}\PY{p}{,} \PY{n}{best\PYZus{}leaf\PYZus{}model}\PY{o}{.}\PY{n}{predict}\PY{p}{(}\PY{n}{X\PYZus{}test}\PY{p}{)}\PY{p}{)}\PY{p}{]}\PY{p}{,}
    \PY{l+s+s1}{\PYZsq{}}\PY{l+s+s1}{F1}\PY{l+s+s1}{\PYZsq{}}\PY{p}{:} \PY{p}{[}\PY{n}{f1\PYZus{}score}\PY{p}{(}\PY{n}{y\PYZus{}test}\PY{p}{,} \PY{n}{tree\PYZus{}gini}\PY{o}{.}\PY{n}{predict}\PY{p}{(}\PY{n}{X\PYZus{}test}\PY{p}{)}\PY{p}{)}\PY{p}{,} \PY{n}{f1\PYZus{}score}\PY{p}{(}\PY{n}{y\PYZus{}test}\PY{p}{,} \PY{n}{tree\PYZus{}pruned}\PY{o}{.}\PY{n}{predict}\PY{p}{(}\PY{n}{X\PYZus{}test}\PY{p}{)}\PY{p}{)}\PY{p}{,} \PY{n}{f1\PYZus{}score}\PY{p}{(}\PY{n}{y\PYZus{}test}\PY{p}{,} \PY{n}{best\PYZus{}leaf\PYZus{}model}\PY{o}{.}\PY{n}{predict}\PY{p}{(}\PY{n}{X\PYZus{}test}\PY{p}{)}\PY{p}{)}\PY{p}{]}\PY{p}{,}
    \PY{l+s+s1}{\PYZsq{}}\PY{l+s+s1}{Profundidade}\PY{l+s+s1}{\PYZsq{}}\PY{p}{:} \PY{p}{[}\PY{n}{tree\PYZus{}gini}\PY{o}{.}\PY{n}{get\PYZus{}depth}\PY{p}{(}\PY{p}{)}\PY{p}{,} \PY{n}{tree\PYZus{}pruned}\PY{o}{.}\PY{n}{get\PYZus{}depth}\PY{p}{(}\PY{p}{)}\PY{p}{,} \PY{n}{best\PYZus{}leaf\PYZus{}model}\PY{o}{.}\PY{n}{get\PYZus{}depth}\PY{p}{(}\PY{p}{)}\PY{p}{]}\PY{p}{,}
    \PY{l+s+s1}{\PYZsq{}}\PY{l+s+s1}{Folhas}\PY{l+s+s1}{\PYZsq{}}\PY{p}{:} \PY{p}{[}\PY{n}{tree\PYZus{}gini}\PY{o}{.}\PY{n}{get\PYZus{}n\PYZus{}leaves}\PY{p}{(}\PY{p}{)}\PY{p}{,} \PY{n}{tree\PYZus{}pruned}\PY{o}{.}\PY{n}{get\PYZus{}n\PYZus{}leaves}\PY{p}{(}\PY{p}{)}\PY{p}{,} \PY{n}{best\PYZus{}leaf\PYZus{}model}\PY{o}{.}\PY{n}{get\PYZus{}n\PYZus{}leaves}\PY{p}{(}\PY{p}{)}\PY{p}{]}
\PY{p}{\PYZcb{}}\PY{p}{)}
\PY{n}{final\PYZus{}comp}
\end{Verbatim}
\end{tcolorbox}

    
    \begin{Verbatim}[commandchars=\\\{\}]
   min\_samples\_leaf  test\_accuracy  test\_precision  test\_recall   test\_f1  \textbackslash{}
0                 1       0.873120        0.459345     0.477001  0.468006   
1                 5       0.886390        0.516476     0.454316  0.483406   
2                10       0.896343        0.574486     0.439824  0.498216   
3                20       0.899882        0.598283     0.439193  0.506541   
4                50       0.898924        0.608216     0.382483  0.469632   
5               100       0.897965        0.613408     0.345936  0.442385   
6               200       0.898186        0.601378     0.385003  0.469458   
7               500       0.897007        0.606502     0.340895  0.436466   

   tree\_depth  leaves  
0          36    2562  
1          33    1411  
2          31     892  
3          26     541  
4          21     269  
5          19     162  
6          15      93  
7          11      43  
    \end{Verbatim}

    
            \begin{tcolorbox}[breakable, size=fbox, boxrule=.5pt, pad at break*=1mm, opacityfill=0]
\prompt{Out}{outcolor}{22}{\boxspacing}
\begin{Verbatim}[commandchars=\\\{\}]
                Modelo  Acurácia  Precisão    Recall        F1  Profundidade  \textbackslash{}
0        Sem restrição  0.873120  0.459345  0.477001  0.468006            36
1          max\_depth=5  0.902094  0.638503  0.376181  0.473434             6
2  min\_samples\_leaf=20  0.899882  0.598283  0.439193  0.506541            26

   Folhas
0    2562
1      61
2     541
\end{Verbatim}
\end{tcolorbox}
        
    \textbf{Resposta:}

\begin{enumerate}
\def\labelenumi{\arabic{enumi}.}
\tightlist
\item
  Escolhi o parâmetro min\_samples\_leaf.
\item
  Ele define o número mínimo de amostras que uma folha precisa ter para
  ser criada, evitando que a árvore crie folhas com 1 ou 2 exemplos para
  decorar ruídos.
\item
  A profundidade reduziu de 36 para 26 e o número de folhas caiu de 2562
  para 541 (com min\_samples\_leaf=20).
\item
  A acurácia de teste subiu de 87,31\% para 89,99\% e o F1-Score subiu
  de 0,4680 para 0,5065 (o melhor F1 entre todos os modelos).
\item
  min\_samples\_leaf (ou a combinação dele com max\_depth), pois permite
  que ramos com muitos dados continuem crescendo enquanto poda ramos sem
  suporte estatístico.
\end{enumerate}


    % Add a bibliography block to the postdoc
    
    
    
\end{document}
